% ---------------------------------------------------------------------------
% Annexe A — Fiches des décisions d'architecture (ADR)
% Source : notes/15-adr-history-principles.md (§1.2, §3, §5.2, §5.3) ;
% docs/adr/ADR-001 à ADR-042 et docs/adr/README.md. Sorties observées avec
% target/debug/reactant (commit e67b10a), NO_COLOR=1, depuis /tmp/manuscrit-adr.
% Le récit chronologique et les fils rouges sont au chapitre 26
% (chap:histoire-doctrine) ; cette annexe est le catalogue de référence, une
% fiche par ADR, citée partout ailleurs par \adr{NN}.
% ---------------------------------------------------------------------------
\chapter{Fiches des décisions d'architecture (ADR)}
\label{chap:fiches-adr}

\begin{objectifs}
  \item Disposer d'une fiche de référence pour chacune des 42 décisions d'architecture (\code{docs/adr/ADR-001} à \code{ADR-042}) : contexte, décision, alternatives refusées, conséquences, statut, ADR liés, fichiers, chapitre du manuscrit qui l'explique.
  \item Savoir retrouver, pour un type ou un mécanisme du code, l'ADR qui le justifie, et inversement.
  \item Distinguer un ADR toujours en vigueur d'un ADR remplacé (\emph{superseded}) ou amendé, et suivre la chaîne jusqu'à la décision actuelle.
\end{objectifs}

\prerequis{aucun chapitre n'est nécessaire pour \emph{consulter} une fiche isolée, mais chaque fiche suppose connu le vocabulaire du \cref{chap:interpretation-abstraite} (treillis, must/may, soundness) ; pour une lecture suivie plutôt qu'un usage de référence, lire d'abord \cref{chap:histoire-doctrine}, qui raconte l'histoire que ces fiches documentent.}

Cette annexe n'est pas un résumé narratif : c'est un catalogue. Chaque
\terme{ADR} (\emph{Architecture Decision Record}) y a une section numérotée
\code{ADR-NN}, avec le label \code{adr:NN} auquel renvoie la macro
\code{\textbackslash adr\{NN\}} employée dans tout le manuscrit. La règle
d'entrée d'un ADR, redite ici, est celle de \file{docs/adr/README.md} : « une
décision que le reste du système doit respecter : un domaine, une relation, un
invariant, une alternative refusée ». Les corrections de précision — une règle
sur-signalait une forme, le moteur a été corrigé, le corpus l'a mesuré — ne
sont pas des ADR et vivent dans \file{docs/precision-log.md}, hors du
périmètre de cette annexe.

Chaque fiche suit huit champs fixes. \textbf{Contexte} pose le problème
concret qui a motivé la décision. \textbf{Décision} résume ce qui a été
tranché, souvent avec une citation de l'ADR original (en anglais : les ADR de
\reactant{} sont rédigés en anglais). \textbf{Refusé} liste les alternatives
examinées et leur raison de rejet — c'est la partie la plus utile à qui
voudrait « nettoyer » le code sans relire l'ADR. \textbf{Conséquences} nomme
les fichiers et types touchés, tels qu'observés à \code{\snapshotCommit}
(certains fichiers annoncés par un ADR n'ont jamais existé, ou ont depuis
disparu : le dire est plus utile que de le taire). \textbf{Statut} reprend le
vocabulaire de \file{docs/adr/README.md} : \emph{Accepted} (décidé, en
vigueur), \emph{Implemented} (décidé et dont l'implémentation complète est
constatée), \emph{Superseded by ADR-NN} (remplacé). \textbf{ADR liés} pointe
vers les amendements et remplacements, dans les deux sens. \textbf{Fichiers}
donne les chemins qui incarnent la décision aujourd'hui. \textbf{Chapitre}
renvoie au chapitre technique qui explique le mécanisme en détail — cette
annexe résume, elle n'explique pas à nouveau.

Quinze ADR sont dits \terme{structurants} : ils fixent une pièce du système
que presque tout le reste présuppose (ADR-001, 002, 003, 009, 010, 012, 014,
015, 017, 018, 021, 022, 027, 041, 042). Leur fiche est plus longue et cite du
code \code{Rust} vérifié dans le dépôt ; les autres tiennent en une dizaine de
lignes. L'ordre suit la numérotation des ADR, qui est aussi, à quelques
regroupements de journée près, l'ordre chronologique (\cref{chap:histoire-doctrine}
donne la frise complète).

% ===========================================================================
\section{ADR-001 --- La sémantique concrète : React-tRace}
\label{adr:1}
\noindent\textit{29 mai 2026. Statut : Accepted.}

\textbf{Contexte.} Une interprétation abstraite calcule une abstraction $\alpha$
d'une sémantique concrète $C$ (\cref{def:concretisation}). Sans $C$ explicite,
la soundness (\cref{def:soundness}) ne peut pas être établie formellement, et
les fonctions de transfert sont écrites « by guesswork ».

\textbf{Décision.} Adopter comme sémantique concrète de référence l'article
\emph{React-tRace} \cite{LeeAhnYi2025} (Lee, Ahn, Yi, OOPSLA 2025), seule
formalisation publique des hooks \code{useState}/\code{useEffect} avec preuve
de conformité empirique. Le modèle \emph{Tree Memory} et sa boucle de rendu
(\code{StepInit} $\to$ \code{StepEffect} $\to$ \code{StepCheck}) correspond
directement à l'itération de point fixe de \reactant{} (\cref{chap:point-fixe-composant}) ;
les règles \code{SttReBind}, \code{CheckEffect}, \code{CheckNoEffect}
définissent exactement les conditions de re-render que l'analyse doit
détecter.

\textbf{Refusé.} Rien d'alternatif n'est débattu : la décision est fondatrice.
Ses limites sont assumées explicitement : React-tRace ne couvre que
\code{useState}/\code{useEffect} \emph{sans} tableau de dépendances ; les
extensions (\code{useMemo}, \code{useCallback}, \code{useRef}, objets)
devaient être spécifiées \og as extensions of React-tRace\fg{}, sans garantie
formelle équivalente.

\textbf{Conséquences.} Trois artefacts étaient annoncés : \file{docs/semantics.md}
spécifiant les extensions ; des fonctions de transfert citant en commentaire
la règle React-tRace correspondante ; des tests de régression contre
l'interpréteur OCaml du papier. \textbf{Aucun des trois n'existe à
\code{\snapshotCommit}} : ni \file{docs/semantics.md}, ni occurrence de
\og React-tRace\fg{} ou \og SttReBind\fg{} dans \file{src/} ou \file{tests/}
(vérifié : \code{ls docs/semantics.md} échoue, \code{grep -r} est vide). La
sémantique concrète reste un repère de conception dans les têtes des
contributeurs, jamais un artefact exécutable ou vérifié mécaniquement.

\textbf{ADR liés.} Aucun (fondateur) ; le README du projet cite toujours le
papier.

\textbf{Fichiers.} Aucun des fichiers annoncés n'a été livré. Le cycle qui en
dérive est \src{src/engine/fixpoint.rs}{704}{704} (\code{analyze_program}).

\textbf{Chapitre.} \cref{chap:react-modele} (le modèle React abstrait) et
\cref{chap:point-fixe-composant} (le cycle qui en dérive).

% ===========================================================================
\section{ADR-002 --- Treillis de stabilité et trois stores}
\label{adr:2}
\noindent\textit{29 mai 2026. Statut : Accepted, raffiné par \adr{17}.}

\begin{react}[\code{Object.is} et la stabilité référentielle]
React compare les dépendances d'un \code{useEffect}/\code{useMemo} et les
props d'un composant mémoïsé par \code{Object.is} (identité de référence pour
un objet, une fonction ou un tableau). Une fonction ou un objet littéral
recréé à chaque rendu est donc \emph{toujours} \og différent\fg{} de sa
version précédente, même si son contenu est identique. C'est cette propriété
— la \emph{stabilité référentielle} — que le domaine doit suivre, pas la
valeur elle-même.
\end{react}

\textbf{Contexte.} Il faut un domaine abstrait pour les références (objets,
tableaux, fonctions) qui capture exactement ce dont les règles ont besoin :
« cette référence est-elle la même à chaque rendu ? ».

\textbf{Décision.} Un treillis à quatre points $\Bot \lleq \{\text{\code{Stable}},
\text{\code{Unstable}}\} \lleq \Top$ ; une table de transfert statique ; trois
stores séparés plutôt qu'un store unifié : \code{StateStore} (soumis au point
fixe), \code{MemoStore} (dérivé), \code{RefStore} (trivial, un \code{ref}
n'entre jamais dans une comparaison de dépendances). Un widening après deux
itérations.

\textbf{Refusé.} Aucune alternative n'est débattue dans cet ADR fondateur.

\textbf{Conséquences.} Fichiers annoncés : \file{stability.rs}, \file{ref_store.rs},
\file{product.rs}. À \code{\snapshotCommit}, \code{Stability} vit dans
\file{src/domains/impls/stability.rs} ; \textbf{\file{src/domains/stores/ref_store.rs}
et \file{src/domains/product.rs} n'existent plus} (vérifié : \code{ls}
échoue sur les deux) — \adr{14} notait déjà que \code{ProductDomain::widen}
était \og a never-wired cartesian product\fg{}. Le treillis à quatre points
lui-même est \textbf{raffiné par \adr{17}} : \code{Unstable} devient
\code{PerRender} (borne must) et gagne \code{Versioned(S)}/\code{VersionedTop}
(borne may) — voir \cref{sec:statevalue-stabilite-stabilite} pour la forme
actuelle.

\textbf{ADR liés.} Raffiné par \adr{17} ; les stores fusionnent plus tard dans
le domaine produit d'\adr{15}.

\textbf{Fichiers.} \file{src/domains/impls/stability.rs} ; \file{src/domains/stores/state_store.rs}.

\textbf{Chapitre.} \cref{chap:treillis-simples} (le treillis, avant
raffinement) et \cref{chap:transfert-stores} (les stores).

% ===========================================================================
\section{ADR-003 --- Une IR dédiée à base de CFG}
\label{adr:3}
\noindent\textit{29 mai 2026. Statut : Accepted.}

\textbf{Contexte.} L'AST produit par \code{oxc} a des dizaines de formes
syntaxiquement différentes pour une même situation sémantique (\code{if}/\code{else}
vs opérateur ternaire, boucle \code{for} vs \code{while}, etc.). Faire porter
les domaines abstraits directement sur cet AST les obligerait à traiter
chaque forme séparément.

\textbf{Décision.} Abaisser (\emph{lowering}) l'AST vers une IR propre,
représentée comme un CFG (\cref{def:cfg}) plutôt qu'un arbre : les retours
précoces et les boucles s'y expriment nativement, les composantes fortement
connexes donnent les points de widening, la dominance (nécessaire pour
détecter les hooks conditionnels, \regle{conditional-hook}) y est standard.
Les hooks sont des nœuds de première classe de l'IR (pas une convention de
nommage retrouvée après coup) ; les composants sont identifiés par un ordre de
priorités : préfixe \code{use} $\Rightarrow$ hook ; JSX retourné $\Rightarrow$
composant ; annotation explicite.

\textbf{Refusé.} Une IR arborescente à la React-tRace : elle exigerait une
transformation en \emph{style de passage de continuation} (CPS) pour
représenter les retours précoces, et ne peut pas représenter de boucle sans
arête arrière — un choix incompatible avec la dominance dont \adr{25} aura
besoin.

\textbf{Conséquences.} Les domaines abstraits ne voient jamais l'AST \code{oxc},
seulement l'IR. \file{docs/ir.md}, annoncé, a disparu ; la table de désucrage
elle-même a dérivé depuis (le diamant de blocs qui désucre \code{\&\&}/\code{||},
conservé par \adr{20} item 1 --- \cref{sec:histoire-doctrine-codification}).

\textbf{ADR liés.} Sa table de désucrage dérive au fil du temps (constat
d'\adr{20}) ; le CFG qu'il pose est raffiné par \adr{25} (fall-through) et
\adr{35} (arête \code{Await}).

\textbf{Fichiers.} \file{src/ir/} (\file{mod.rs}, \file{cfg.rs}, \file{expr.rs},
\file{stmt.rs}, \file{hooks.rs}) ; \file{src/lowering/} (abaissement en une
passe).

\textbf{Chapitre.} \cref{chap:ir}.

% ===========================================================================
\section{ADR-004 --- \code{render_cfg} et \code{effect_cfg} séparés}
\label{adr:4}
\noindent\textit{29 mai 2026. Statut : Accepted.}

\textbf{Contexte.} Le rendu et les effets ont des sémantiques distinctes
vis-à-vis du \code{StateStore} : un rendu lit l'état, un effet peut l'écrire et
en déclencher un nouveau. Les confondre dans un seul graphe rendrait cette
distinction implicite.

\textbf{Décision.} \code{ComponentIR\{ render_cfg, hooks: Vec<HookEntry> \}} ;
chaque effet porte son propre \code{body_cfg}. Le cycle d'analyse est
rendu $\to$ effets $\to$ join $\to$ test de convergence $\to$ widening
(\cref{chap:point-fixe-composant}).

\textbf{Refusé.} Un méta-CFG unifié, avec des arêtes arrière
rendu\,$\to$\,effet\,$\to$\,vérification. Plus expressif, mais jugé
\og hard to maintain\fg{} et rendant implicite la séparation render/effet, qui
est un invariant de React : autant la rendre structurelle. Les bugs visés
(boucles infinies, état incohérent) sont des propriétés de l'état au point
fixe, pas des propriétés de chemins dans un graphe combiné.

\textbf{Conséquences.} Aucun changement depuis : c'est toujours la forme du
code à \code{\snapshotCommit}.

\textbf{ADR liés.} Aucun amendement direct ; présuppose \adr{3} (l'IR CFG) et
conditionne \adr{9} (traversée des callbacks pendant les effets).

\textbf{Fichiers.} \file{src/ir/component.rs} (\code{ComponentIR}) ;
\file{src/engine/fixpoint.rs} (le cycle rendu/effets).

\textbf{Chapitre.} \cref{chap:ir} (la structure) et \cref{chap:point-fixe-composant}
(le cycle).

% ===========================================================================
\section{ADR-005 --- Portée intra-procédurale et registre de hooks}
\label{adr:5}
\noindent\textit{29 mai 2026. Statut : Accepted, dépassé de fait.}

\textbf{Contexte.} Un composant appelle des hooks personnalisés et des hooks de
bibliothèques (\code{react-router}, \code{TanStack Query}\ldots). L'analyse ne
peut pas tous les inliner.

\textbf{Décision.} Une première phase intra-procédurale : un hook inconnu
vaut $\Top$ (\code{Unknown}). Un registre \code{HookModel} en couches
(hooks natifs, modules de bibliothèques connus, configuration utilisateur
\code{reactant.toml}, repli \code{Unknown}).

\textbf{Refusé.} Rien de débattu ; la portée intra-procédurale est un choix
d'étape, explicitement provisoire.

\textbf{Conséquences.} \file{src/registry/user_config.rs} et \code{reactant.toml}
étaient annoncés — \textbf{jamais créés} (vérifié : \code{src/registry/user_config.rs}
n'existe pas). À \code{\snapshotCommit}, \file{src/registry/} contient
\file{keyed.rs}, \file{mod.rs}, \file{summary.rs}. Le registre en couches est
dépassé de fait par l'inlining inter-composants (\adr{12}), l'inlining
cross-file (\adr{13}) et le \code{SummaryRegistry} — des résumés packagés pour
des bibliothèques précises (TanStack, React Router, \code{next/navigation})
plutôt qu'une configuration utilisateur.

\textbf{ADR liés.} Dépassé par \adr{12}, \adr{13} et le \code{SummaryRegistry}.

\textbf{Fichiers.} \file{src/registry/keyed.rs}, \file{src/registry/summary.rs}.

\textbf{Chapitre.} \cref{chap:frontend-detection}.

% ===========================================================================
\section{ADR-006 --- Les règles en post-passe}
\label{adr:6}
\noindent\textit{29 mai 2026. Statut : Accepted.}

\textbf{Contexte.} Une règle peut être écrite \emph{inline}, couplée aux
fonctions de transfert du point fixe, ou en \emph{post-passe} sur le résultat
convergé.

\textbf{Décision.} Post-passe pure sur \code{AnalysisResult}. Seule exception
« inline » tolérée : enregistrer \code{widened_labels} pendant le point fixe
lui-même (une information que la post-passe ne pourrait pas reconstituer).
\og The engine […] doesn't know about rules\fg{}.

\textbf{Refusé.} Le couplage inline généralisé : il romprait la frontière
moteur/règles et forcerait chaque nouvelle règle à comprendre l'algorithme de
point fixe.

\textbf{Conséquences.} La signature \code{(\&AnalysisResult) -> Vec<Warning>}
est remplacée par \code{check(\&RuleCtx)} (\adr{21} §4) ; la frontière est
durcie plus tard par \adr{42} (les règles ne marchent plus aucune syntaxe,
seulement des relations).

\textbf{ADR liés.} Sa signature est remplacée par \adr{21} ; sa frontière est
durcie par \adr{42}.

\textbf{Fichiers.} \src{src/rules/registry.rs}{254}{254} (\code{RuleRegistry::check_component}).

\textbf{Chapitre.} \cref{chap:api-regles}.

% ===========================================================================
\section{ADR-007 --- Requêtes inter-domaines}
\label{adr:7}
\noindent\textit{2 juin 2026. Statut : Accepted.}

\textbf{Contexte.} Un domaine (par exemple un futur \code{SetterEffect}) doit
parfois lire l'état d'un autre domaine (la \code{Stability} d'une référence).
Référence prise chez MOPSA, dont le \emph{Manager} résout ce problème avec des
GADT.

\textbf{Décision (A, implémentée).} Un trait \code{QueryContext} en
\code{\&dyn} (object-safe) ; \code{AnalysisCtx<D>} regroupe
\code{(state, memo, heap, query)} ; trois implémentations
(\code{NullCtx}, \code{FixpointCtx}, \code{AnalysisQueryCtx}).
\textbf{Décision (B, différée).} Des types marqueurs \code{Queryable<Q>} en
Rust stable, la spécialisation (\code{\#[feature(specialization)]}, RFC
\#31844) étant instable.

\textbf{Refusé.} Le \emph{Manager} générique à la MOPSA : les GADT sont
absents de Rust stable et la spécialisation est instable — la décision B ne
peut pas être prise tout de suite.

\textbf{Conséquences.} \adr{15} refuse ensuite de généraliser en \og query
pool\fg{} $n$-aire ; \adr{21} réalise le \og typed Manager later\fg{} promis
ici, sous une autre forme (la surface de requêtes typée des règles).

\textbf{ADR liés.} La généralisation en pool est refusée par \adr{15} ; le
\og Manager typé plus tard\fg{} est réalisé par \adr{21}.

\textbf{Fichiers.} \file{src/domains/mod.rs} (\code{AnalysisCtx}, trait
\code{QueryContext}).

\textbf{Chapitre.} \cref{chap:transfert-stores}.

% ===========================================================================
\section{ADR-008 --- \code{StateValue} : enum plat et \code{TypedStateStore}}
\label{adr:8}
\noindent\textit{2 juin 2026. Statut : Superseded by \adr{15}.}

\textbf{Contexte.} La \code{Stability} seule ne distingue pas
\code{setState(count + 1)} (croît sans borne) de \code{setState(42)} (converge
immédiatement) : il faut une vraie valeur abstraite pour la partie état.

\textbf{Décision.} Un enum plat \code{StateValue \{ Bottom, Null, Undefined,
Number(Interval), \ldots, Top \}} ; des sous-stores typés par sorte inférée
statiquement ; puis (4 juin) un indice \code{useState<T>} pour contourner le
problème que \code{join(Null, Number) = Top} fait perdre toute borne.

\textbf{Refusé (limite résiduelle assumée).} \code{useState(null)} sans
annotation de type reste un faux négatif : sans indice, la sorte \code{Null}
et la sorte \code{Number} restent incomparables, le \code{join} retombe à
$\Top$.

\textbf{Conséquences.} \textbf{Superseded by \adr{15}} le 14 juillet : les
trois mécanismes (l'enum plat, \code{TypedStateStore}, l'indice
\code{useState<T>}) sont \emph{supprimés} en bloc ; \code{Interval},
\code{BoolVal}, \code{StrConst} survivent comme composantes du domaine
produit.

\textbf{ADR liés.} Superseded by \adr{15} (le seul ADR entièrement remplacé
plutôt qu'amendé, \cref{sec:histoire-doctrine-corpus}).

\textbf{Fichiers.} Aucun : \code{TypedStateStore} n'existe plus à
\code{\snapshotCommit}.

\textbf{Chapitre.} \cref{chap:statevalue-stabilite} (pour la forme qui l'a
remplacé).

% ===========================================================================
\section{ADR-009 --- Traversée sémantique des callbacks}
\label{adr:9}
\noindent\textit{2 juin 2026, complétée le 4 juin. Statut : Accepted --- complete.}

\textbf{Contexte.} \code{fetch().then(u => setUser(u))} est invisible pour un
point fixe qui ne regarde que les instructions du corps direct. Mais
descendre uniformément dans \emph{tout} callback armerait un faux positif : un
\code{addEventListener} de clic verrait son handler exécuté à chaque
itération du point fixe, faisant croître un état qui ne croît en réalité qu'au
clic.

\textbf{Décision.} Classer chaque appelé : \code{TriggerClass} distingue
\code{Setter} (le setter lui-même, jamais descendu), \code{InCycle}
(conséquence directe du rendu ou de l'effet courant : \code{map}, \code{forEach},
\code{.then}, \code{setTimeout}\ldots --- descendu), \code{Subscription}
(abonnement externe, \emph{non} descendu) et \code{Unknown} (appelé non
reconnu, \emph{non} descendu). Les \emph{handlers} deviennent des points
d'entrée séparés de l'analyse (\code{HookEntry::Handler},
\code{extract_handlers}, \code{extract_subscriptions}) ; un widening provoqué
par un handler est exclu de \code{widened_labels}.

\rustsnippet{src/domains/interp/callbacks.rs}{6}{23}{\code{TriggerClass} --- la classification qui décide si un appelé est descendu}

\textbf{Refusé --- choix assumé non sound.} \code{Unknown} → \code{skip} :
\og for a linter, false positives are more costly than false negatives.
Descending an unknown callee would be the \emph{sound} choice […] We accept
the FN\fg{}. C'est un choix \emph{antérieur} à \file{CLAUDE.md} (juin, contre
juillet) et il contredit aujourd'hui l'invariant \og faux négatifs
interdits\fg{}.

\begin{piege}
Ce choix a survécu, mais pas partout : la \textbf{marche des relations}
(\adr{34} §5) a pris la décision inverse --- elle descend tout argument
fonction d'un appel inconnu, à $\Top$, \og because an unknown callee may
invoke it\fg{}. Deux traversées du même code, deux politiques différentes sur
les appelés inconnus : c'est le résidu tracé par les issues \issue{19} (FN,
appelés inconnus sans \code{Loc}) et \issue{12} (deux interpréteurs de force
différente). Ne pas supposer qu'une correction faite d'un côté vaut pour
l'autre.
\end{piege}

\textbf{Conséquences.} \file{src/domains/interp/callbacks.rs} porte
\code{classify_callee} et \code{TriggerClass}. \file{src/domains/interp/interpreter.rs}
porte \code{exec_callbacks_depth} (appelé dans \code{exec_full_stmt}) et
\code{exec_body}/\code{exec_body_impl}, la descente \og effets de bord
seulement, valeur de retour ignorée\fg{}. L'API publique (\code{Transfer},
règles, \code{AnalysisResult}) est inchangée pour ce premier incrément.

\textbf{ADR liés.} Le choix \code{Unknown} → \code{skip} est inversé pour la
marche des relations par \adr{34} §5.

\textbf{Fichiers.} \file{src/domains/interp/callbacks.rs} ; \file{src/domains/interp/interpreter.rs}.

\textbf{Chapitre.} \cref{chap:point-fixe-composant}.

% ===========================================================================
\section{ADR-010 --- Modèle de tas par site d'allocation}
\label{adr:10}
\noindent\textit{3 juin 2026. Statut : Accepted --- complete.}

\textbf{Contexte.} Deux formes échappaient à l'analyse : un callback stocké
dans une variable puis passé ailleurs (\code{const cb = () => \ldots; setTimeout(cb)}),
et un appel direct à une fonction locale (\code{load()} définie plus haut dans
le composant).

\textbf{Décision.} Un identifiant \code{ExprId} sur chaque nœud allouant
(littéral objet, tableau, fonction) ; un tas \code{Heap: ExprId -> HeapValue}
monotone ; l'environnement abstrait porte deux cartes séparées,
\code{stabs} (valeurs) et \code{locs} (localisations), car une carte unique
\code{Val | Loc} s'écrasait sur elle-même dès qu'une variable changeait de
rôle. Une profondeur d'inlining plafonnée :

\rustsnippet{src/domains/interp/interpreter.rs}{21}{21}{la profondeur d'inlining des appels locaux}

\textbf{Refusé --- limite connue assumée.} Au-delà de cette profondeur, ou en
cas de récursion, l'appel n'est pas descendu : un diagnostic \Info{}
\regle{analysis-limit} le signale plutôt que de le passer sous silence.

\textbf{Conséquences.} \adr{23} §3 révélera plus tard un défaut de conception
laissé par cet ADR : \code{locs} est monotone et n'est \emph{jamais invalidé}
à la réaffectation d'une variable --- un raffinement futur, pas un
non-changement décidé.

\textbf{ADR liés.} Défaut révélé (non corrigé ici) par \adr{23} §3.

\begin{sloppypar}
\textbf{Fichiers.} \file{src/ir/types.rs} (\code{ExprId}) ;
\file{src/domains/stores/abstract_env.rs} (\code{locs}) ;
\file{src/domains/stores/heap.rs} ;
\file{src/engine/cfg_analyzer.rs} et \file{src/engine/fixpoint.rs} (le tas
créé une fois, partagé entre passes de rendu et d'effets).
\end{sloppypar}

\textbf{Chapitre.} \cref{chap:ir} (\code{ExprId}) et \cref{chap:transfert-stores}
(le tas).

% ===========================================================================
\section{ADR-011 --- \code{SourceRange} et notes}
\label{adr:11}
\noindent\textit{3 juin 2026. Statut : Accepted, §Note superseded by \adr{19}.}

\textbf{Contexte.} Sans position, un diagnostic ne peut pas être affiché à
côté du code fautif ; et un seul type de note limitait l'analyse à un seul
participant nommé par finding.

\textbf{Décision.} \code{SourceRange \{ line, col \}} ; une table
\code{line_starts} construite en $O(n)$, offrant une conversion en
$O(\log n)$ par recherche dichotomique. \code{Note \{ message, hook_label,
range \}} ; \code{Option<SourceRange>} pour ne pas casser les tests d'IR
écrits à la main sans position réelle.

\textbf{Refusé.} Rien de débattu au-delà de la structure retenue.

\textbf{Conséquences.} La structure \code{Note} et sa limitation \og pas de
fichier\fg{} (un seul projet ne pouvait alors désigner qu'un fichier
implicite) sont \textbf{superseded by \adr{19}}, qui ajoute \code{FileId}
après l'arrivée de l'inlining cross-file (\adr{13}) : une note pouvait sinon
pointer dans le mauvais fichier.

\textbf{ADR liés.} §Note superseded by \adr{19}.

\textbf{Fichiers.} \src{src/ir/source_range.rs}{36}{43}{}.

\textbf{Chapitre.} \cref{chap:ir}.

% ===========================================================================
\section{ADR-012 --- Analyse inter-composants}
\label{adr:12}
\noindent\textit{4 juin 2026. Statut : Accepted.}

\textbf{Contexte.} Un setter passé en \emph{prop} à un enfant, puis appelé
là-bas, ferme une boucle qui traverse la frontière entre deux composants ---
invisible à une analyse strictement intra-composant.

\textbf{Décision.} Inlining \textbf{descendant} (\emph{top-down)} : le parent
est analysé d'abord, jamais de résumés remontants (\emph{bottom-up}), \og
without justified gain at the current stage\fg{}. Un cache indexé par égalité
abstraite des props (double \code{leq}, dans les deux sens) évite de
ré-analyser un enfant avec des props équivalentes. \code{ComponentSetter}
(l'analogue du \code{Set\_clos} de React-tRace) ; \code{SharedStateStore} pour
l'état partagé entre composants inlinés ; un \code{RootDetector} modulaire
(\code{--all-roots}, \code{--entry}) pour décider par où commencer. La
récursion entre composants donne $\Top$ (même politique que MOPSA pour la
récursion de fonctions).

\textbf{Refusé.} Les résumés remontants (bottom-up) : pas de gain démontré à
ce stade du projet.

\textbf{Conséquences.} Limites reconnues : les imports non résolus, les
composants dynamiques (\code{const C = cond ? A : B}, \issue{63}, plus tard
\code{wontfix}, \cref{sec:histoire-doctrine-methode}). \code{src/engine/fixpoint.rs}
reçoit \code{SharedStateStore} en paramètre et propage les appels
\code{ComponentSetter} ; \code{src/engine/} gagne \code{ComponentRegistry},
\code{RootDetector}, \code{ProgramAnalysisResult} ; les règles migrent vers
cette dernière.

\textbf{ADR liés.} L'inlining bottom-up refusé ici reste refusé jusqu'à
\code{\snapshotCommit}.

\begin{sloppypar}
\textbf{Fichiers.} \file{src/domains/stores/} (\code{SharedStateStore}) ;
\file{src/engine/} (\code{ComponentRegistry}, \code{RootDetector}) ;
\file{src/resolver/mod.rs}, l.~439 (\code{analyze_lowered}).
\end{sloppypar}

\textbf{Chapitre.} \cref{chap:inter-composants}.

% ===========================================================================
\section{ADR-013 --- Analyse cross-file}
\label{adr:13}
\noindent\textit{5 juin 2026. Statut : Accepted --- phases 1--4.}

\textbf{Contexte.} Trois problèmes à la fois : des collisions de noms
(\code{Page()} défini dans plusieurs fichiers) ; une découverte de fichiers
manuelle ; des utilitaires opaques qui masquent un setter (\code{doOrNot(setX(...))}
devient un faux positif si \code{doOrNot} n'est jamais regardé).

\textbf{Décision.} Une clé \code{(PathBuf, String)} plutôt qu'un nom seul ;
deux traits, \code{FileDiscoverer} et \code{ImportResolver} ; une entrée par
répertoire ; un graphe de \textbf{symboles} (pas de fichiers) trié
topologiquement pour ordonner l'analyse ; \code{FunctionIR} et l'inlining des
fonctions utilitaires ; une analyse \emph{eager} (tout le graphe, pas de
paresse) ; un import non résolu vaut $\Top$ plus un \Info{} --- \og FPs
possible, FNs forbidden\fg{}, déjà la politique de soundness du projet.

\textbf{Refusé.} Rien d'alternatif débattu pour la structure retenue ; les
limites sont assumées et listées, pas contournées : alias \code{tsconfig},
chaînes de ré-export, \code{node_modules} (jamais abaissé, \issue{51},
\code{wontfix}), inlining seulement en position d'instruction (\issue{52}).

\textbf{Conséquences.} \file{src/resolver/mod.rs} porte \code{lower_files}
(\src{src/resolver/mod.rs}{237}{237}{}) et \code{lower_files_with}
(\src{src/resolver/mod.rs}{247}{247}{}).

\textbf{ADR liés.} Le périmètre \code{node_modules} refusé ici est confirmé
\code{wontfix} par l'issue \issue{51} (\cref{sec:histoire-doctrine-methode}).

\textbf{Fichiers.} \file{src/resolver/mod.rs} ; \file{src/ir/function.rs}
(\code{FunctionIR}).

\textbf{Chapitre.} \cref{chap:inter-composants}.

% ===========================================================================
\section{ADR-014 --- Widening à seuils ; narrowing abandonné}
\label{adr:14}
\noindent\textit{27--28 juin 2026. Statut : Accepted (partie 1) ; partie 2 superseded.}

\textbf{Contexte.} Le widening standard d'un intervalle saute directement à
$\pm\infty$ : \code{count} borné par une garde \code{count < 10} converge vers
$[0, \infty)$, perdant l'information \code{count} $\in [0, 10]$ que la garde
donnait pourtant.

\textbf{Décision (partie 1, implémentée).} Récolter des \emph{seuils} dans les
littéraux des gardes et des initialisations rencontrées dans le CFG ;
\code{widen_to(\&self, other, \&[f64])} qui élargit vers le seuil le plus
proche au lieu de $\pm\infty$, sans changer la signature de \code{widen}
elle-même ; la même technique s'applique à l'arête arrière interne
d'\code{analyze_cfg}.

\textbf{Refusé --- partie 2, abandonnée le lendemain de sa proposition.} Un
\emph{narrowing} descendant, pour affiner après convergence. Révision
interne : « narrowing and threshold widening both recover only concrete
literal bounds » --- le narrowing n'aurait rien apporté que le widening à
seuils n'apportait déjà, et une phase descendante ne peut de toute façon pas
faire redescendre un store dont les \code{join} sont monotones (une valeur ne
peut que croître au sens du treillis). L'ADR exigeait des tests de propriété
pour cette partie 2 : ils n'ont jamais été écrits, faute d'objet.

\textbf{Conséquences.} \file{src/domains/impls/interval.rs} (\code{widen_to});
\file{src/domains/impls/state_value.rs} route les seuils vers la composante
\code{num} ; \file{src/engine/fixpoint.rs} calcule les seuils avant la boucle
externe et les transmet à \code{widen_to} ; \file{src/engine/analysis_result.rs} :
\code{widened_labels: HashSet} devient plus tard une \code{HashMap} d'un enum
\code{WidenOutcome}.

\textbf{ADR liés.} Aucun amendement ultérieur.

\textbf{Fichiers.} \file{src/domains/impls/interval.rs} ; \file{src/engine/fixpoint.rs} ;
\file{src/engine/cfg_analyzer.rs}.

\textbf{Chapitre.} \cref{chap:cfg-analyzer-dominance}.

% ===========================================================================
\section{ADR-015 --- Domaine produit sur sortes disjointes}
\label{adr:15}
\noindent\textit{14 juillet 2026. Statut : Implemented ; supersedes \adr{8}.}

\textbf{Contexte.} Après \adr{8}, chaque \code{join} entre deux sortes JS
différentes (un \code{null} et un \code{number}) retombait à $\Top$ : trois
rustines empilées (\code{TypedStateStore}, inférence de type, indice
\code{useState<T>}) tentaient de compenser sans traiter la cause.

\textbf{Décision.} \og disjoint sum $\Rightarrow$ union = product\fg{} : les
sortes primitives JS sont mutuellement exclusives (une valeur n'est jamais à
la fois un nombre et une chaîne), donc l'union disjonctive des sortes dégénère
en un \emph{produit} --- un emplacement indépendant par sorte, chacun $\Bot$
quand cette sorte est impossible.

\rustsnippet{src/domains/impls/state_value.rs}{16}{44}{\code{StateValue} --- le domaine produit sur sortes disjointes}

Les opérations (\code{join}, \code{meet}, \code{widen}) sont ponctuelles,
sorte par sorte : un \code{join} entre un état \code{null} et un état
\code{number} garde \emph{les deux} sortes (\code{number|null}) au lieu de
s'effondrer sur $\Top$. C'est précisément ce qui rend détectable un compteur
initialisé à \code{null} : la coercition JS \code{ToNumber(null) = 0} est
câblée dans les opérateurs binaires, et un narrowing de nullabilité et de
véracité s'ajoute sur les gardes (\code{if (x)}, \code{x ?? d}, \code{x !== null}).

\begin{exemple}{Le FN d'ADR-008 devient une détection}{fn-null-counter-adr015}
\begin{lstlisting}[language=TypeScript,caption={Compteur initialisé à \code{null}},label={lst:null-counter-adr015}]
import { useState, useEffect } from "react";

function Counter() {
  const [n, setN] = useState(null);
  useEffect(() => {
    setN(n + 1);
  }, [n]);
  return <span>{n}</span>;
}
\end{lstlisting}
Sous \adr{8}, \code{n} valait $\Top$ dès l'initialisation (\code{join(Null, Number) = Top}),
et un $\Top$ ne peut jamais prouver de croissance : silence. Sous \adr{15},
\code{n} porte la sorte \code{null} \emph{et} la sorte \code{num = [0,0]} ;
\code{n + 1} coerce le \code{null} en $0$ et produit \code{num = [1,1]}, qui
croît à chaque itération du point fixe. Sortie observée (\code{--no-color},
\code{\snapshotCommit}) :
\begin{sortie}
  Counter  (2 hooks)  null-counter.tsx
    warn   infinite-loop  [hook:0]  (line 5:2)  this effect keeps pushing state `n` (its deps do not provably gate it, so the effect can re-run every render) to new values on every run. Potential infinite render loop
       (2 trace step(s), rerun with --trace)

⚠  1 warning(s) across 1 file(s).
\end{sortie}
Le niveau observé est \Warning{}, pas \Error{} : les deps \code{[n]} ne sont
pas \emph{prouvablement} une garde (\cref{chap:infinite-loop} explique la
distinction entre le bras \emph{déps} et le bras \emph{churn} d'\regle{infinite-loop}).
\end{exemple}

\textbf{Refusé.} Un opérateur de réduction inter-sortes, ou un \og query
pool\fg{} $n$-aire à la MOPSA (repris d'\adr{7}) : \og the slots describe
disjoint kinds\fg{}, rien à réduire puisque rien ne se recouvre.

\textbf{Conséquences.} \file{src/domains/impls/state_value.rs} (la struct
produit, les opérations ponctuelles, \code{to_stability} \og motion-wins\fg{}) ;
\file{src/domains/impls/setter_val.rs} (\code{SetterVal}, l'ancienne variante
\code{ComponentSetter} d'\adr{12} §7 devenue une sorte à treillis plat) ;
\file{src/ir/hooks.rs} : \code{HookEntry::State} perd son \code{type_hint}.
Un ancien FN devient une détection épinglée
(\src{tests/narrowing.rs}{221}{221}{\code{null_init_without_hint_unbounded_is_flagged}}) ;
un non-FP est épinglé symétriquement
(\src{tests/narrowing.rs}{326}{326}{\code{nullable_fetch_pattern_no_false_positive}}).
Commit \code{c32e1b0} : 40 fichiers, +1\,674/$-$1\,198 lignes.

\textbf{ADR liés.} Supersedes \adr{8} ; \code{KindMask} (\file{state_value.rs},
\og the product's kind enumeration lives in exactly one place\fg{}) est un
apport ultérieur d'\adr{20} (item D4, \cref{sec:histoire-doctrine-codification}) ;
\adr{20} item 10 interdit d'y câbler \code{TSType}.

\textbf{Fichiers.} \file{src/domains/impls/state_value.rs} ; \file{src/domains/impls/setter_val.rs} ;
\file{src/ir/hooks.rs}.

\textbf{Chapitre.} \cref{chap:statevalue-stabilite}.

% ===========================================================================
\section{ADR-016 --- CLI, sortie JSON, détection de projet Vite}
\label{adr:16}
\noindent\textit{15 juillet 2026. Statut : Implemented.}

\textbf{Contexte.} Le binaire dupliquait le pipeline \emph{sans} le résolveur
d'imports (un faux négatif silencieux sur tout code inatteignable depuis un
point d'entrée mal détecté) ; il n'y avait pas de sortie machine ; les projets
Vite n'étaient pas analysables (structure de répertoires non reconnue).

\textbf{Décision.} Des sous-commandes (\code{check} par défaut, \code{rules},
\code{explain}) ; un pipeline \emph{unique} partagé
(\code{lower_files}, \code{analyze_lowered}) entre le binaire et les tests ;
les métadonnées des règles indexées par nom de \emph{diagnostic} (onze
structures de règles émettaient treize noms distincts à l'époque) ; un DTO
JSON côté binaire ; un contrat de codes de sortie 0/1/2 (succès, warning-seul,
error) ; \code{ProjectKind::\{Plain, Vite\}} ; support \code{tsconfig} JSONC
avec \code{extends} et \code{references}.

\textbf{Refusé.} Évaluer \code{vite.config.*} pour résoudre les alias :
\og a best-effort regex would risk \emph{wrong} resolutions, i.e. silent
FNs\fg{} --- le risque de mal résoudre un chemin (et donc de rater du code)
l'emporte sur le confort d'une détection plus riche.

\textbf{Conséquences.} Le schéma JSON est passé en v2 depuis (\issue{129},
\cref{sec:histoire-doctrine-methode}).

\textbf{ADR liés.} Le schéma JSON évolue par la suite (v2, \issue{129}) sans
ADR dédié (correction de précision de l'ergonomie, pas de la soundness).

\textbf{Fichiers.} \src{src/driver/mod.rs}{106}{106}{} (\code{run_check}) ;
\file{src/driver/json.rs}, \file{src/driver/human.rs}.

\textbf{Chapitre.} \cref{chap:projet-cli-driver}.

% ===========================================================================
\section{ADR-017 --- Stabilité versionnée}
\label{adr:17}
\noindent\textit{15 juillet 2026. Statut : Implemented.}

\textbf{Contexte.} Un couple faux positif/faux négatif noué ensemble : un
état objet stable en pratique (recréé seulement au \code{setState}) était lu
comme \code{Unstable}, produisant des FP \og always-unstable-deps\fg{} sur les
context providers du corpus \emph{memos} ; et \og any FP fix that silences
that warning without a replacement signal silently drops a true infinite
loop\fg{} --- l'\code{ObjChurn} correspondant restait invisible.

\textbf{Décision.} Introduire des bornes \emph{may} (\og ne change qu'aux
\code{set} \fg{}) et \emph{must} (\og change à chaque rendu\fg{}), au lieu du
seul point \code{Unstable}.

\rustsnippet{src/domains/impls/stability.rs}{35}{64}{\code{Stability} --- le treillis versionné : \code{Versioned}/\code{VersionedTop} (may) et \code{PerRender} (must)}

La conversion \code{Versioned} $\to$ lecture d'état se fait \textbf{côté
lecture}, en un seul endroit du transfert (le bras \code{Expr::StateVal(l)}) :
le store garde une vue \emph{événementielle} de l'état (une valeur par
\code{set}), l'évaluation en tire une vue \emph{inter-rendus}. Un nouveau bras
\og churn\fg{} d'\regle{infinite-loop} est stratifié \Error{}/\Warning{} en
plus du bras \og deps\fg{} historique.

\textbf{Refusé.} Un compteur d'événements pour faire diverger artificiellement
le store à chaque \code{set} : \og a non-standard hack\fg{}, alors que la
certitude doit venir de la structure des dépendances, pas d'un artifice de
comptage. Deux points intermédiaires du produit complet (\emph{may}-seul,
\emph{must}-seul) ont aussi été envisagés puis écartés au profit du couple
complet.

\textbf{Couplage de soundness.} Le \emph{gating} par \code{Versioned} n'est
sound \emph{qu'avec} le bras churn actif : sans lui, un état
\og seulement \emph{may}-changeant\fg{} laisserait passer un cas qui change en
réalité à chaque rendu.

\textbf{Conséquences.} Cinq FP \og memos\fg{} supprimés ; les bugs de la
classe \code{ObjChurn} passent de \og noyés dans le bruit\fg{} à \Error{}
propre. \code{Stability} perd \code{Copy} (elle porte désormais un ensemble) :
\og the compile errors force an exhaustive audit of every consumer ---
desirable for a semantics change\fg{}. Commit \code{f31acae} : 26 fichiers,
+2\,364 lignes.

\textbf{ADR liés.} Raffine \adr{2} ; couplé au graphe de churn d'\adr{18} pour
la soundness du gating.

\textbf{Fichiers.} \file{src/domains/impls/stability.rs}.

\textbf{Chapitre.} \cref{chap:statevalue-stabilite} et \cref{chap:infinite-loop}
(le bras churn).

% ===========================================================================
\section{ADR-018 --- Graphe de churn multi-effets}
\label{adr:18}
\noindent\textit{16 juillet 2026. Statut : Implemented ; amendé par \adr{42}.}

\textbf{Contexte.} Une boucle répartie sur \emph{deux} effets --- l'effet A
écrit un slot lu par l'effet B, qui écrit à son tour un slot lu par A --- est
invisible à un bras \regle{infinite-loop} qui ne regarde qu'un effet à la
fois.

\textbf{Décision.} Un graphe sur les slots qualifiés \code{(component, label)} ;
une arête \og un changement de \code{x} relance un effet qui stocke une
référence fraîche dans \code{y}\fg{} ; une arête \emph{must} exige une
dépendance exacte \emph{et} une fraîcheur (\code{Fresh}) sur tous les chemins ;
une auto-arête pour les effets sans tableau de deps ; une \og convergence
kill\fg{} conditionnée à un \textbf{écrivain unique} du slot ; l'algorithme de
Tarjan en deux passes pour les composantes fortement connexes. Le graphe est
construit une fois par programme (résolvant un problème de performance
quadratique, \issue{86}) ; les cycles inter-composants sont plafonnés à
\Warning{}.

\textbf{Refusé.} Rien de débattu au-delà de la construction retenue ; la
condition d'écrivain unique est une simplification \emph{assumée}, pas
débattue ici --- elle sera remplacée plus tard.

\textbf{Conséquences.} Cycles objets sur deux ou $N$ effets : \Info{} $\to$
\Error{} avec le chemin du cycle. Écritures fraîches sans tableau de deps
(l'oubli classique) : silence $\to$ \Error{}. Churn objet inter-composants
(sous condition \code{Versioned}) : silence $\to$ \Warning{}. Corpus (4
dépôts) : inchangé, aucun nouveau FP.

\textbf{ADR liés.} Amendé par \adr{42} §4 et §6 : la construction migre dans
le moteur, et la condition d'écrivain unique est remplacée par la preuve
multi-sites de l'issue \issue{154}.

\textbf{Fichiers.} À \code{\snapshotCommit}, le graphe vit dans
\file{src/engine/churn.rs} (déplacé depuis \file{rules/helpers/} par \adr{42} ---
voir cette fiche).

\textbf{Chapitre.} \cref{chap:churn}.

% ===========================================================================
\section{ADR-019 --- Chaînes de témoins typées}
\label{adr:19}
\noindent\textit{17 juillet 2026. Statut : Implemented ; supersedes \adr{11} §Note.}

\textbf{Contexte.} Les notes de diagnostic étaient du texte libre, sans
fichier associé, et la provenance d'un fait (pourquoi le moteur croit-il
cela ?) était jetée dès qu'elle avait servi.

\textbf{Décision.} \code{FileId}/\code{FileTable} ; la provenance est
enregistrée \og au point de connaissance\fg{} (\code{widen_trace},
\code{inline_origins}) plutôt que reconstruite après coup ; un vocabulaire
\emph{fermé} \code{Step} (neuf variantes à l'origine : \code{Binding},
\code{Resolve}, \code{Call}, \code{Write}, \code{Read}, \code{Branch},
\code{Handler}, \code{CycleEdge}, \code{Widen}), \textbf{sans} variante
\code{Text(String)} ; un rendu centralisé, une bibliothèque de témoins
partagée entre règles ; des bornes (un niveau de résolution, huit étapes
maximum).

\textbf{Refusé.} Un domaine de provenance complet, à la MOPSA :
\og touches every transfer function and multiplies memory for witness depth
no rule needs\fg{}. Une variante \code{Step::Text(String)} : \og a free-text
escape hatch would erode the vocabulary back to prose within months\fg{}.

\textbf{Conséquences.} À \code{\snapshotCommit}, l'enum \code{Step} compte
cinq variantes de plus (\code{Mutate}, \code{Capture}, \code{InitOnce},
\code{Forward}, \code{Rerender}) : le vocabulaire a crû par la voie
sanctionnée --- l'extension de l'enum --- jamais par une variante texte libre.
Commit \code{146a86a} : 55 fichiers, +1\,737/$-$322 lignes ; tous les
littéraux \code{SourceRange} du dépôt sont touchés.

\textbf{ADR liés.} Supersedes \adr{11} §Note.

\textbf{Fichiers.} \file{src/rules/api/witness.rs} (\code{enum Step}, ligne 86).

\textbf{Chapitre.} \cref{chap:api-regles}.

% ===========================================================================
\section{ADR-020 --- Dette technique : onze non-changements délibérés}
\label{adr:20}
\noindent\textit{23 juillet 2026. Statut : Accepted.}

\textbf{Contexte.} Une campagne d'audit sur dix-huit \og workarounds\fg{}
suspects et trente-six constats d'architecture a trouvé que beaucoup de
\og duplications\fg{} apparentes étaient en réalité des variations porteuses
de soundness --- les retirer aurait introduit un faux négatif.

\textbf{Décision.} Une liste de changements réellement appliqués (splice
unifié avec $\alpha$-renommage, \code{recompute_memo} corrigé pour utiliser
les vrais stores, vocabulaire churn hissé dans un module partagé, une
sentinelle opaque typée, la macro \code{flat_lattice!}, \code{KindMask}\ldots)
et \textbf{onze non-changements}, chacun avec sa phrase de justification
(reproduits en détail au \cref{sec:histoire-doctrine-codification}, tableau~\ref{tab:histoire-adr020}) :
garder le diamant \code{\&\&}/\code{||} plutôt qu'un \code{LogicalOp} plat ;
garder les deux bras du churn séparés ; garder \code{may_written_slots}
syntaxique ; garder la projection \code{to_stability} ; ne pas généraliser en
\code{BoundedPowerset<T,N>} ; ne pas partager une résolution de composant
unique entre règles ; garder séparé le chemin \og position de retour\fg{} ;
garder le seed de tas par site d'appel ; garder \code{resolve_setter_aliases}
comme passe séparée ; ne pas câbler \code{TSType} dans le domaine ; garder les
noms temporaires par \emph{offset}.

\textbf{Refusé.} Chacun des onze changements listés ci-dessus, précisément
parce qu'il coûtait un faux négatif identifiable sur un programme concret.

\textbf{Conséquences.} \file{CLAUDE.md} cite désormais cet ADR comme lecture
obligatoire avant toute nouvelle tentative de \og nettoyage\fg{}. La leçon est
répétée mot pour mot dans le dépôt : \og map before ``fixing''; never
introduce a false negative to remove a duplication\fg{}.

\textbf{ADR liés.} Réinvoqué explicitement par \adr{29} (item 2, refus de
fusionner les deux bras churn) et \adr{31} (item 3, \code{may_written_slots}).

\textbf{Fichiers.} Aucun fichier propre : c'est un ADR de non-changement, sa
trace est dans \file{docs/adr/ADR-020-tech-debt-cleanup-decisions.md} lui-même.

\textbf{Chapitre.} \cref{chap:histoire-doctrine}, \cref{sec:histoire-doctrine-codification}
(table complète des onze non-changements).

% ===========================================================================
\section{ADR-021 --- Surface de requêtes typée}
\label{adr:21}
\noindent\textit{24 juillet 2026. Statut : Accepted --- implemented.}

\textbf{Contexte.} Trois défauts distincts : un FN réel dans un \emph{helper}
partagé (\code{is_unstable} n'était pas exactement le complément de
\code{is_stable}) ; \code{Severity::Error} attachable en pratique à
n'importe quel booléen, sans lien démontrable avec une preuve must ; le
\og must-forward\fg{} (propager une preuve certaine à travers un appel)
dupliqué dans chaque règle qui en avait besoin. Étude préalable : cinq
\emph{frontends} de règles prototypés, avec la conclusion \og the frontend is
not the soundness lever --- the query surface is\fg{}.

\textbf{Décision.} La polarité devient un \textbf{type} plutôt qu'une
convention :

\rustsnippet{src/rules/api/query.rs}{111}{129}{\code{MustResult<T>} --- le verdict must à trois valeurs, et \code{May<T>}, sans chemin vers une \Error{}}

\code{Diagnostic::error} n'est atteignable que depuis un \code{Certified} :

\rustsnippet{src/rules/api/query.rs}{74}{84}{\code{Certified<E>} --- le jeton scellé : constructeur privé au module}

Des primitives \og graines\fg{} produisent les premiers \code{Certified} ;
\code{RuleCtx} regroupe l'accès aux relations. Migration \og big-bang\fg{} des
quatorze règles existantes en un seul changement : \og a transition window
would leave the FN vector open\fg{} --- laisser cohabiter ancienne et nouvelle
API, même temporairement, aurait laissé un chemin non vérifié vers un faux
négatif.

\begin{piege}
La première version du scellement n'était \emph{pas} étanche : la visibilité
Rust étant descendante (un module voit ses parents publics), \code{Diagnostic::new}
restait appelable depuis les sous-modules de règles, et un champ
\code{pub severity} permettait de \emph{forger} une \code{Error} sans passer
par un \code{Certified}. Le \og hardening round\fg{} a déplacé
\code{Diagnostic} dans un module feuille dédié, \file{src/rules/api/diagnostic.rs}
(commit \code{eb5cb93}), pour qu'aucun module frère ne puisse plus le
construire directement. Retenir la leçon : un scellement de type ne protège
que ce que la visibilité du module protège réellement.
\end{piege}

\textbf{Refusé.} Une fenêtre de transition où l'ancienne signature
\code{(\&AnalysisResult) -> Vec<Warning>} resterait disponible en parallèle de
la nouvelle : elle laisserait ouvert le vecteur de FN le temps de la
migration.

\textbf{Conséquences.} \textbf{Positif} : aucune \code{Error}-sur-\emph{may} ni
aucun abandon silencieux de $\Top$ ne survit à la compilation ; le
must-forward est partagé et à trois valeurs ; la surface de requêtes devient
l'ancre stable sur laquelle le futur \emph{frontend} de règles se lie ; le FN
vérifié est corrigé et épinglé par un test. \textbf{Coût} : les quatorze
règles et la signature du trait migrent en un seul changement ; la rigueur de
relecture exigée sur les annotations de polarité augmente.

\textbf{ADR liés.} Remplace la signature d'\adr{6} ; un amendement du 1{\ier}
septembre introduit \code{ProgramCache}, repris ensuite par \adr{42} §5.

\textbf{Fichiers.} \file{src/rules/api/query.rs} ; \file{src/rules/api/diagnostic.rs}.

\textbf{Chapitre.} \cref{chap:api-regles}.

% ===========================================================================
\section{ADR-022 --- Frontends de règles, packs déclaratifs, WASM}
\label{adr:22}
\noindent\textit{25 juillet 2026. Statut : Accepted, §7 superseded by \adr{23} §5, §4 amendé par \adr{41} §4.}

\textbf{Contexte.} Les cinq questions laissées ouvertes par l'étude
préalable d'\adr{21} : comment une équipe écrit-elle sa propre règle sans
toucher au moteur ni en perdre la soundness ?

\textbf{Décision.} Principe normatif : \og reactant does \emph{semantics
only}\fg{} --- une règle personnalisée ne lit jamais l'AST ni le texte source,
seulement des relations sémantiques déjà résolues par le moteur (les
\og ancres\fg{} Tier-A). Une règle est une ancre plus une navigation typée
dans les relations qu'elle expose. La sévérité est \code{pin} $\meet$ \code{polarity}
--- le minimum entre un plafond choisi par l'auteur de la règle et la
polarité du fait lui-même --- évaluée \textbf{par finding}, ce qui donne la
stratification \Error{}/\Warning{} \og gratuitement\fg{} (sans code
supplémentaire par règle). Paramètres restreints à une position de constante
feuille (jamais une expression arbitraire). Espace de noms \code{pack/rule} ;
documentation obligatoire par règle ; distribution \textbf{WASM-only} sur
npm ; le cœur Rust revalide systématiquement tout pack chargé, même compilé
--- \og the host JS is not a trust boundary\fg{}.

\textbf{Refusé.} Le rejet statique des conflits \code{pin}/\code{polarity} :
il aurait interdit la stratification par finding, qui est précisément ce que
le $\meet$ rend possible. Les callbacks JS bruts avec autorité
d'\Error{} : un callback JS ne peut jamais être certifié par construction.

\textbf{Conséquences.} \textbf{Positif} : les règles personnalisées héritent
de la soundness du moteur par construction ; un seul mécanisme de sévérité
($\meet$) sert auteurs de règles, consommateurs et règles natives ; le
schéma JSON dérive de la surface \code{api/}, donc la croissance du
vocabulaire Tier-C se propage gratuitement ; WASM-only garantit un
comportement identique partout. \textbf{Coût} : il faut construire une couche
\og entité-arête\fg{} par-dessus \code{api/} (les primitives prennent des
types IR bruts --- \code{\&CFG}, \code{\&Expr} --- non JSON-appelables
directement) ; les messages d'erreur du validateur doivent être assez bons
pour nourrir une boucle d'auteur assistée par un LLM.

\textbf{ADR liés.} §7 (une syntaxe déclarative en Starlark) superseded by
\adr{23} §5 ; §4 (la portée des options de règles) amendé par \adr{41} §4.

\textbf{Fichiers.} \file{src/rules/api/} (surface) ; distribution WASM (\cref{chap:distribution}).

\textbf{Chapitre.} \cref{chap:regles-declaratives} et \cref{chap:distribution}.

% ===========================================================================
\section{ADR-023 --- Croissance du vocabulaire Tier-A}
\label{adr:23}
\noindent\textit{26 juillet 2026. Statut : Accepted ; supersedes \adr{22} §7.}

\textbf{Contexte.} Mesuré sur le catalogue : seules 3 règles sur 21 étaient
alors exprimables avec les seules ancres et primitives Tier-A disponibles.

\textbf{Décision.} Croître le vocabulaire par \textbf{entités} --- des
positions dans des relations déjà résolues --- plutôt que par des gardes ad
hoc. Un verdict doit se lire \emph{au point de programme de son entité}
(sinon c'est un nouveau fait non prouvé, donc un FN potentiel). Première
entité : \code{args}, avec une primitive \code{returns_verdict} --- amendée le
lendemain pour être calculée \textbf{pendant} le point fixe, car le contexte
nécessaire ne survit pas à la post-passe.

\textbf{Refusé.} Le quantificateur universel $\forall$ sur les éléments d'un
tableau : un $\Top$ y plierait en \og viole\fg{}, et une vacuité sur une arête
tronquée serait un faux positif silencieux. \textbf{Amendement du
1{\ier} septembre} : la condition est levée par un bit \code{exact} porté par
les littéraux \code{ArrayLit} ; la primitive \code{every} est livrée, mais
seulement au positif, sans autorité d'\Error{}. Starlark comme langage de
règles déclaratives est rejeté : il ne compile pas nativement dans l'arbre de
dépendances du projet, qui a déjà \code{oxc} sous la main pour parser du
JS/TS et le compiler en JSON.

\textbf{Conséquences.} Le catalogue passe de 5/21 à 8/22 (le dénominateur
change avec l'entité \code{args}), puis continue de croître ADR après ADR
jusqu'à 21/22 (\adr{34}).

\textbf{ADR liés.} Supersedes \adr{22} §7 (Starlark) ; le refus du
quantificateur $\forall$ est partiellement levé par l'amendement du 1{\ier}
septembre.

\textbf{Fichiers.} \file{src/rules/api/} (entité \code{args}) ; le calcul de
\code{returns_verdict} vit dans le point fixe
(\src{src/engine/fixpoint.rs}{130}{130}{\code{analyze_component_impl}}).

\textbf{Chapitre.} \cref{chap:regles-declaratives}.

% ===========================================================================
\section{ADR-024 --- Attribution des findings inlinés}
\label{adr:24}
\noindent\textit{26 juillet 2026. Statut : Accepted.}

\textbf{Contexte.} Sur un échantillon de packs, 12 findings sur 27 (44\,\%)
s'affichaient sous le fichier du composant \emph{consommateur} avec la ligne
du hook \emph{inliné} --- une attribution qui pointait le mauvais fichier pour
un contributeur qui voudrait corriger le hook lui-même.

\textbf{Décision.} La ligne primaire d'un finding rend l'\textbf{origine}
(le fichier du hook) quand elle diffère du fichier consommateur ; en JSON, le
champ \code{file} devient le fichier de l'ancre plutôt que celui du composant
affiché.

\textbf{Refusé.} Dédupliquer les findings entre consommateurs d'un même hook
inliné (\code{useStep(1)} déclenchant \regle{infinite-loop} et
\code{useStep(0)} déclenchant \regle{redundant-set-state} : ce sont des faits
\emph{incomparables}, fusionner perdrait de l'information). Un groupement
d'affichage différé : repris plus tard, autrement, par l'issue \issue{129},
qui groupe seulement à l'\emph{affichage} tout en gardant une ligne par
composant en JSON.

\textbf{Conséquences.} Aucune régression du nombre de findings ; l'attribution
change, pas le compte.

\textbf{ADR liés.} Le groupement d'affichage refusé ici revient sous une
forme différente avec l'issue \issue{129}.

\textbf{Fichiers.} \file{src/driver/json.rs}, \file{src/driver/human.rs}.

\textbf{Chapitre.} \cref{chap:api-regles}.

% ===========================================================================
\section{ADR-025 --- Fall-through = \code{return undefined}}
\label{adr:25}
\noindent\textit{29 juillet 2026. Statut : Accepted.}

\textbf{Contexte.} \code{Terminator::Unreachable} confondait trois situations
distinctes : la fin normale d'un corps de fonction sans \code{return}
explicite, un \code{throw}, et un \code{break} orphelin. Le \emph{splice}
(greffe d'un CFG dans un autre, \cref{chap:splice-inlining}) laissait alors le
bloc de jointure sans prédécesseur valide. Mesuré : 208 composants
\og sectionnés\fg{} (l'environnement de sortie, \code{exit_env}, incomplet ---
un faux négatif).

\textbf{Décision.} Un corps qui \og tombe\fg{} à sa fin sans \code{return}
explicite est scellé \code{Return(Lit(Unit))} --- un \code{return undefined}
implicite, exactement la sémantique JS. \code{throw} garde
\code{Unreachable}. \code{ExitDominance} devient le seul propriétaire de
l'ensemble des sorties \textbf{atteignables} d'un CFG.

\rustsnippet{src/ir/cfg.rs}{12}{30}{\code{Terminator} --- \code{Return} sans \code{span}, \code{Unreachable} réservé à l'arrêt de contrôle}

\begin{piege}
Le champ \code{span} de \code{Terminator::Return} est \code{None} par
construction (contrairement à \code{Branch}, qui en porte un) : \og nothing
needed the position of a \code{return} until a hook could be extracted from
one\fg{}. Un hook atteint seulement par un chemin qui finit sur un
\code{return} produira donc un finding \emph{sans numéro de ligne} ---
visible, mais non localisé. C'est un choix assumé (issue \issue{140} encore
ouverte), pas un oubli. Le commentaire du code chiffre ce coût à \og a 40-site
IR change\fg{} ; six semaines plus tard, refusant le même ajout pour une
autre raison (nommer l'appelant plutôt que l'appelé lors d'un \emph{splice}),
\adr{39} le chiffre à \og a hundred construction sites\fg{}. Les deux
estimations, écrites à deux moments par le même changement refusé, ne se
recoupent pas --- signe qu'il s'agit d'ordres de grandeur à main levée, pas
d'une mesure, et qu'aucune des deux ne doit être citée comme un chiffre
d'audit.
\end{piege}

\textbf{Refusé --- mesuré.} Relier tout \code{Unreachable} au bloc de
jointure : cela aurait fait naître une \Error{} \regle{conditional-hook} sur
l'idiome courant \og garde puis \code{throw}\fg{} :

\begin{exemple}{L'idiome guard-throw, non signalé}{ex-guard-throw}
\begin{lstlisting}[language=TypeScript,caption={\code{useCart} : garde par exception avant les hooks}]
function useCart(items: unknown[]) {
  if (!items) {
    throw new Error("items required");
  }
  const [open, setOpen] = useState(false);
  return { open, setOpen };
}
\end{lstlisting}
Sortie observée (\code{\snapshotCommit}) : \emph{aucun} finding.
\begin{sortie}
   1 clean component(s) hidden, rerun with --show-clean

✓  1 file(s) no issues found.
\end{sortie}
Relier le \code{throw} au join aurait rendu \code{useState} atteignable par
deux chemins de hauteur différente dans l'arbre des appels de hooks, un motif
que \regle{conditional-hook} confondrait avec un hook réellement conditionnel.
\end{exemple}

\textbf{Conséquences.} Douze findings révélés, zéro perdu ; les composants
\og sectionnés\fg{} tombent de 208 à 3. Commit \code{34ce48b} : 7 fichiers,
+395 lignes.

\textbf{ADR liés.} Aucun amendement direct ; \adr{39} révise le calcul des
positions manquantes en aval.

\textbf{Fichiers.} \file{src/ir/cfg.rs} (\code{Terminator}) ;
\file{src/engine/dominance.rs} (\code{ExitDominance}).

\textbf{Chapitre.} \cref{chap:lowering-cfg} et \cref{chap:cfg-analyzer-dominance}.

% ===========================================================================
\section{ADR-026 --- Projets Next.js}
\label{adr:26}
\noindent\textit{27 août 2026. Statut : Implemented.}

\textbf{Contexte.} Les projets Next.js n'étaient analysés \og only by
accident\fg{} (aucune reconnaissance de leur structure). Question de fond :
que signifie un \emph{Server Component} pour un analyseur bâti sur le
rendu/l'état/les effets, des notions qui n'existent pas côté serveur ?

\textbf{Décision.} \textbf{Ne pas sauter} les modules serveur --- sauter sur
la base d'arêtes incomplètes du graphe d'imports serait un FN. \code{ModuleFacts \{
directives, imports \}} : les directives (\code{"use client"}, \code{"use
server"}) et les imports sont enregistrés \emph{sans être interprétés} comme
une preuve de domaine ; \code{reachable_from} calcule l'atteignabilité dans ce
graphe. \code{ProjectKind::NextJs} testé \emph{avant} \code{Vite} (un projet
Next.js a aussi un \code{vite.config}, rarement, mais la priorité doit être
sans ambiguïté). \code{baseUrl} comme résolveur de chemins additionnel. La
règle \regle{server-component-hook} est plafonnée \Warning{} : les faits
qu'elle utilise (convention de nommage, graphe d'imports) sont hors du
domaine sémantique de l'analyse. Des résumés pour \code{next/navigation}.

\textbf{Refusé.} Une must-primitive certifiant qu'un module est un Server
Component : \og would dress a path convention as a domain proof\fg{} --- la
convention \code{"use client"} n'est pas un fait sémantique prouvé par
interprétation abstraite, seulement une convention de fichier. Supprimer les
autres findings dans un module serveur : cela masquerait de vrais bugs
d'état/effets dans le code qui s'exécute malgré tout côté serveur au premier
rendu.

\textbf{Conséquences.} Zéro FP sur cinq corpus Next.js testés ; 46 \Info{}
\regle{unknown-hook} supprimés (les hooks de Next.js désormais reconnus) ;
quatre findings révélés par la nouvelle règle.

\textbf{ADR liés.} Aucun amendement.

\textbf{Fichiers.} \file{src/resolver/mod.rs} (\code{ProjectKind::NextJs}) ;
la règle \regle{server-component-hook}.

\textbf{Chapitre.} \cref{chap:regles-rendu-rsc} et \cref{chap:projet-cli-driver}.

% ===========================================================================
\section{ADR-027 --- Relation slot-writer, résumés de phase, provenance}
\label{adr:27}
\noindent\textit{1{\ier} septembre 2026. Statut : Accepted, §2 amendé par \adr{34} et \adr{35}.}

\textbf{Contexte.} Des règles de type \og wrapper enforcement\fg{}
(\og l'état ne s'écrit que via \code{putState}\fg{}, une politique d'équipe)
étaient inexprimables : rien ne disait \emph{où}, ni \emph{dans quelle
phase}, un slot était écrit. Un audit du code contre les ADR a motivé une
relation dédiée.

\textbf{Décision.} Une ligne d'écrivain à deux colonnes distinctes :
\code{region} (\emph{exacte} --- \code{Render}, \code{Effect}, etc.) et
\code{phase} (\emph{may}, $\Top$ pour un \code{FnLit} imbriqué non encore
résolu).

\rustsnippet{src/engine/setters.rs}{681}{701}{\code{WriterPhase} --- la phase d'exécution d'une écriture, un verdict may}

La relation est calculée \emph{à convergence} du point fixe et
\textbf{déplacée des règles vers le moteur} (première étape du fil rouge
\og relations = produits du moteur\fg{}, \cref{sec:histoire-doctrine-relations}).
Les résumés de phase des appelés sont \emph{whitelistés} avant que le
vocabulaire ne soit publié --- affiner $\Top$ après coup changerait des
findings déjà livrés, ce qui violerait la stabilité du catalogue. Résolution
d'imports des fonctions utilitaires \og fail-closed\fg{} (import non résolu
$\to$ $\Top$). Provenance enregistrée au site unique du splice.
\code{must_direct_write} donne l'autorité d'\Error{} aux règles de politique
d'équipe. Catalogue re-basé à 22 entrées.

\textbf{Refusé.} Classer tout \code{FnLit} imbriqué comme \code{Deferred}
plutôt que $\Top$ : \code{arr.forEach(x => setX(x))} à l'intérieur d'un effet
s'exécute \emph{synchrone} dans la phase de l'effet, pas \og en différé\fg{} ---
classer systématiquement en \code{Deferred} sous-approximerait.

\textbf{Conséquences.} La note d'\adr{23} §4 sur l'ordonnancement de
\code{any_of}/\code{eval_guard} est déclarée périmée. Le test-cliquet
\code{catalogue_has_21_entries} est remplacé par le pin à 22 entrées quand §6
atterrit. \file{docs/custom-rules.md} et les compétences \code{skills/reactant-rules/}
gagnent un test anti-dérive.

\textbf{ADR liés.} §2 (résumés de phase, \code{await}) amendé par \adr{34} et
\adr{35}, qui lèvent les points laissés en suspens.

\textbf{Fichiers.} \file{src/engine/setters.rs} (\code{WriterPhase}, tête de
\code{SlotWriter}).

\textbf{Chapitre.} \cref{chap:relations}.

% ===========================================================================
\section{ADR-028 --- \code{writers} par site, \code{Updater}, \code{same_tick}}
\label{adr:28}
\noindent\textit{1{\ier} septembre 2026. Statut : Accepted.}

\textbf{Contexte.} \regle{stale-update} (deux écritures non fonctionnelles du
même slot dans un même tick d'événement, où seule la dernière compte) était
inexprimable : la relation d'écriture précédente effondrait plusieurs appels
en une seule ligne.

\textbf{Décision.} \textbf{Renverser} l'effondrement documenté par la ligne
précédente : une ligne par \emph{site d'appel}, pas par région.

\rustsnippet{src/engine/setters.rs}{728}{745}{\code{SlotWriter} --- une ligne par site, et pourquoi le renversement est monotone}

Une \textbf{unique} colonne \code{Updater} d'où chaque consommateur dérive son
propre verdict (fonctionnel ou non, pureté du corps de l'\code{Updater})
plutôt que plusieurs colonnes spécialisées. Un booléen \code{same_tick} par
ligne, jamais un pli agrégé --- la relation \og deux écritures dans le même
tick\fg{} reste \emph{may} par construction.

\textbf{Refusé.} Deux colonnes spécifiques au lieu d'un \code{Updater}
partagé : aurait dupliqué la même information sous deux formes qui pourraient
diverger. Une forme \emph{niée} de \code{same_tick} (\og pas dans le même
tick\fg{}) : la relation est \emph{may}, sa négation ne serait pas
\emph{must} et n'apporterait rien.

\textbf{Conséquences.} Le catalogue Tier-A passe de 14/22 à 16/22 ;
\issue{61} reste ouverte (elle ne couvre que la moitié asynchrone du
problème).

\textbf{ADR liés.} L'alternative à deux colonnes, refusée ici, est
explicitement rappelée par \adr{27} §4 (\og one central relation, never a
second bespoke one\fg{}).

\textbf{Fichiers.} \file{src/engine/setters.rs}.

\textbf{Chapitre.} \cref{chap:relations}.

% ===========================================================================
\section{ADR-029 --- Ancre \code{churn_cycles}}
\label{adr:29}
\noindent\textit{1{\ier} septembre 2026. Statut : Accepted, §1 amendé par \adr{42}.}

\textbf{Contexte.} \regle{cross-component-effect-cycle} était bloquée par
l'issue \issue{68} : le catalogue Tier-A n'acceptait alors qu'une seule ancre
par règle, et un cycle de churn est une propriété du \emph{programme}, pas
d'un composant.

\textbf{Décision.} Projeter la relation programme (le graphe de churn) sur le
composant \og ancré\fg{} : une ligne par arête de cycle portée par l'un de ses
effets. Les colonnes sont booléennes et \emph{exactes} --- le caractère
\emph{may} du fait reste dans le graphe lui-même, pas dans la ligne projetée.
Une arête sans position source (\emph{span}) ne produit pas de ligne.

\textbf{Refusé.} Un statut \og couvert nativement, ne compte pas\fg{} pour
sortir cette classe de bug du calcul de couverture Tier-A : la mesure porte
sur le \emph{vocabulaire} disponible aux règles personnalisées, pas sur ce que
les règles natives couvrent déjà par ailleurs. Fusionner les deux bras du
churn (repris et refusé par \adr{20} item 2).

\textbf{Conséquences.} Le catalogue passe de 16/22 à 17/22 ; l'issue
\issue{68} cesse d'être bloquante pour cette classe de règles.

\textbf{ADR liés.} §1 (la construction du graphe) amendé par \adr{42}, qui la
déplace dans le moteur.

\textbf{Fichiers.} L'ancre \code{churn_cycles} lit \file{src/engine/churn.rs}
(après le déplacement d'\adr{42} ; côté règles à l'époque de cet ADR).

\textbf{Chapitre.} \cref{chap:churn}.

% ===========================================================================
\section{ADR-030 --- Lignes render-setter qualifiées par propriétaire}
\label{adr:30}
\noindent\textit{1{\ier} septembre 2026. Statut : Accepted.}

\textbf{Contexte.} \regle{setter-called-in-child-render} était traitée comme
une jointure ad hoc entre deux relations, au lieu de lire une relation dédiée.

\textbf{Décision.} Des lignes \og étrangères\fg{} dérivées de la résolution
\code{cross_component_setters} déjà existante, avec une colonne \code{owner}
(le composant \emph{propriétaire} du slot, pas celui qui appelle le setter).
L'énumération des cas s'élargit \emph{sur la garde} qui sélectionne les lignes
(\code{slot_ownership}), jamais sur la \emph{sorte} de la ligne elle-même ---
élargir la sorte changerait la forme des findings déjà livrés par les packs.
Le nom du slot étranger est affiché dans le composant \textbf{propriétaire},
pas dans celui qui le lit.

\textbf{Refusé.} Élargir la sorte de ligne inconditionnellement : changerait
les findings des packs déjà distribués. Restreindre le scan au seul bloc du
site d'appel : supprimerait des findings dans les cas où l'environnement
abstrait est imprécis --- un coût en soundness pour un gain de bruit
seulement.

\textbf{Conséquences.} Le catalogue passe de 17/22 à 18/22. Aveu explicite :
l'attribution du propriétaire n'était pas exacte au moment de cet ADR
(issue \issue{119}, ouverte puis corrigée par le commit \code{ce3b080}).

\textbf{ADR liés.} Aucun amendement direct au-delà de la correction
\code{ce3b080}.

\textbf{Fichiers.} \file{src/engine/setters.rs} (colonne \code{owner}, voir
aussi \adr{42} §2, où elle est formalisée sur \code{SlotWriter}).

\textbf{Chapitre.} \cref{chap:relations}.

% ===========================================================================
\section{ADR-031 --- Relation \code{slot_seeds}}
\label{adr:31}
\noindent\textit{1{\ier} septembre 2026. Statut : Accepted, §2 amendé par l'issue \issue{121} et \adr{33}.}

\textbf{Contexte.} \regle{state-mirrors-prop-without-sync} était vue comme une
jointure ad hoc entre une prop et un slot d'état, sans relation propre pour
la relier.

\textbf{Décision.} \og A fold promoted to the engine\fg{} : les fonctions
utilitaires de \file{frozen_initial_state.rs} migrent dans
\file{engine/seeds.rs}. La moitié \og effet\fg{} de la relation est
\textbf{dérivée} de \relation{slot_writers} (aucun second parcours du CFG) ;
la moitié \og render\fg{} lit la \textbf{phase prouvée} du moteur, pas la
région lexicale du code source. L'échappement d'une prop hors de sa synchro
initiale est représenté comme une \emph{colonne} de la relation, jamais comme
un verdict de synchronisation déjà tranché.

\textbf{Refusé.} Plier l'échappement dans un champ \code{synced} binaire :
effacerait la distinction précise dont dépend le palier \Error{} de la règle.

\begin{piege}
La première mesure de cet ADR affirmait une sortie native \og inchangée sur
14 corpus\fg{}. C'était \textbf{faux} : un seul run avait été fait par côté de
la comparaison, alors que l'analyse n'est pas rigoureusement déterministe à
ce stade (issue \issue{120}, corrigée ensuite par \adr{33}). Recompté
honnêtement : 12 corpus identiques, \code{twenty} $-1$ (un FP retiré, une
amélioration), \code{mantine} $-1$ (un \emph{vrai} positif perdu par accident,
restauré ensuite par \adr{33}). La leçon : une mesure \og inchangé\fg{} sur un
pipeline non déterministe doit être répétée avant d'être publiée.
\end{piege}

\textbf{Conséquences.} Amendement de l'issue \issue{121} : un filtre
\code{effect_triggered(phase)} affine la moitié effet.

\textbf{ADR liés.} §2 amendé par l'issue \issue{121} ; le déterminisme requis
pour la mesure est livré par \adr{33}.

\textbf{Fichiers.} \file{src/engine/seeds.rs}.

\textbf{Chapitre.} \cref{chap:relations}.

% ===========================================================================
\section{ADR-032 --- Relation \code{context_consumers}}
\label{adr:32}
\noindent\textit{1{\ier} septembre 2026. Statut : Accepted.}

\textbf{Contexte.} \regle{consumer-without-provider} était bloquée sur
l'issue \issue{28} (comment modéliser \code{useContext} sans domaine dédié à
part entière ?).

\textbf{Décision.} Une post-passe purement diagnostique : le verdict de cette
règle est une \textbf{absence} (aucun provider trouvé), donc toute la
conception se ramène à une porte d'ascendance fiable. \emph{Gate 1} :
ascendance complète prouvée dans l'arbre d'éléments (issue \issue{110}).
\emph{Gate 2} : une passe syntaxique de complétion pour les composants non
atteints par l'ascendance sémantique. Appariement sur la cellule canonique
\code{ContextId} (issue \issue{109}). Cas toléré : le provider d'un contexte
compte lui-même comme un \og hit\fg{} pour son propre \code{useContext}
éventuel. Le verdict est nommé \code{none-on-analyzed-paths}, délibérément
pas \code{no-provider} --- le nom dit la limite de l'analyse, pas une absence
prouvée dans l'absolu.

\textbf{Refusé.} Rien d'alternatif débattu pour la structure des deux portes.

\textbf{Conséquences.} Le catalogue passe de 19/22 à 20/22.

\textbf{ADR liés.} Aucun amendement direct.

\textbf{Fichiers.} La relation \relation{context_consumers} (moteur, résolution
d'ascendance dans l'arbre d'éléments).

\textbf{Chapitre.} \cref{chap:relations}.

% ===========================================================================
\section{ADR-033 --- Bit d'exactitude de la chasse aux liaisons}
\label{adr:33}
\noindent\textit{2 septembre 2026. Statut : Accepted.}

\textbf{Contexte.} La sortie de l'analyseur n'était pas rigoureusement
déterministe (issue \issue{120}) : \code{normalize_to_prop} pouvait répondre
un chemin de prop différent à chaque exécution sur le même code.

\textbf{Décision.} L'ensemble \code{seen} de la chasse aux liaisons (qui évite
les cycles de résolution) est cloné \emph{par branche} plutôt que partagé ---
clé sur les racines, récursion indépendante par chemin. La sélection à
travers un objet littéral partiel réutilise \code{members_after_last_spread},
déjà partagée avec l'interpréteur. \code{NormPath \{ path, exact \}} : un
chemin \emph{élargi} (par exemple parce qu'une des branches passait par une
propriété dynamique) ne peut plus soutenir une revendication must.

\textbf{Refusé.} Corriger uniquement §1 (le clonage par branche) sans ajouter
le bit \code{exact} de §3 : aurait seulement \emph{augmenté} la suppression de
findings sans garantir qu'aucun nouveau chemin élargi ne serait pris pour un
must. §3 est \og la moitié qui garde l'autre moitié sound\fg{}.

\textbf{Conséquences.} Le corpus \code{mantine} gagne $+1$ (le vrai positif
perdu par \adr{31} est restauré) ; 13 corpus restent identiques. Effet de
bord positif : \code{collect_component_setter_vars} devient déterministe au
passage.

\textbf{ADR liés.} Corrige le défaut de déterminisme révélé par \adr{31}.

\textbf{Fichiers.} La chasse aux liaisons (\code{normalize_to_prop},
\code{NormPath}) dans le moteur des relations.

\textbf{Chapitre.} \cref{chap:relations}.

% ===========================================================================
\section{ADR-034 --- Relation d'enregistrement, une table de registrars}
\label{adr:34}
\noindent\textit{2 septembre 2026. Statut : Accepted.}

\textbf{Contexte.} Trois lecteurs différents entretenaient chacun sa propre
liste (\code{REGISTRARS}, \code{DEFERRING_GLOBALS}, \code{DEFERRING_METHODS}) ;
le résumé de phase promis par \adr{27} §2 n'avait jamais été écrit.

\textbf{Décision.} Une seule table \code{\&[Registrar]}, avec une colonne
\code{timing} :

\rustsnippet{src/engine/registrations.rs}{40}{57}{\code{Timing} --- \og le contrat, pas le nom\fg{}}

\code{Handler} est réservé à \code{addEventListener} \emph{seulement} (contrat
DOM : aucun appel synchrone possible) ; \code{Unknown} couvre \code{subscribe}/\code{on}/\code{addListener},
dont le contrat n'exclut pas un appel synchrone (RxJS \code{BehaviorSubject}
émet parfois immédiatement au nouvel abonné). Un appariement tri-valué entre
enregistrement et \emph{teardown} (\code{Paired} seul est une revendication
exploitable) --- amendement de l'issue \issue{124} : trois formes de teardown
distinguées. Une écriture faite depuis une classe \code{Handler} ne ferme pas
un cycle de churn (issue \issue{93}). Un \emph{teardown} n'est jamais
compté comme un rappel du callback lui-même.

\begin{react}[le contrat DOM de \code{addEventListener}]
Le DOM ne dispatch jamais un événement de façon synchrone depuis l'appel
\code{addEventListener} lui-même : le callback enregistré ne peut provablement
pas s'exécuter \emph{pendant} l'appel qui l'enregistre. C'est ce contrat, et
lui seul, qui autorise \code{Timing::Handler} ; une bibliothèque d'observables
comme RxJS ne l'offre pas.
\end{react}

\textbf{Refusé.} \code{Handler} pour \code{subscribe}/\code{on} : RxJS peut
émettre \emph{synchroniquement} à l'abonnement, donc aucun contrat ne
garantit l'absence d'exécution immédiate --- classer ces registrars en
\code{Handler} serait le seul sens où \og rétrécir $\Top$\fg{} pourrait perdre
un vrai positif sans preuve.

\textbf{Conséquences.} Le catalogue passe de 20/22 à \textbf{21/22}, décrit
dans l'ADR lui-même comme \og the honest ceiling\fg{}. Amendement de l'issue
\issue{116} : une ancre publique \code{registrations} est ajoutée --- la
décision \code{wontfix} de l'issue \issue{42} (heuristique de nom
d'émetteur pour \regle{stale-closure}) couvre désormais une surface plus
large grâce à cette ancre.

\textbf{ADR liés.} Livre le résumé de phase promis par \adr{27} §2 ; amendé
par les issues \issue{116} et \issue{124}.

\textbf{Fichiers.} \file{src/engine/registrations.rs}.

\textbf{Chapitre.} \cref{chap:relations}.

% ===========================================================================
\section{ADR-035 --- Frontière de phase \code{await}}
\label{adr:35}
\noindent\textit{2 septembre 2026. Statut : Accepted, §1 amendé par \adr{36} §6.}

\textbf{Contexte.} \code{AwaitExpression} était effacée au lowering, sans
trace dans le CFG. La supposition \og \code{sync_phase} : lexis = execution,
provably\fg{} --- c'est-à-dire que la position lexicale d'une instruction
suffit à connaître sa phase d'exécution --- devenait \textbf{fausse} après un
\code{await} : le code qui suit dans le texte s'exécute en réalité sur un tour
ultérieur de la boucle d'événements.

\textbf{Décision.} Scinder le bloc du CFG par une \textbf{arête}
\code{EdgeKind::Await} (une variante d'arête, pas un champ ajouté à
\code{BasicBlock}) :

\rustsnippet{src/ir/cfg.rs}{32}{45}{\code{EdgeKind} --- l'arête \code{Await}, contrôle inchangé, phase rompue}

\code{post_await_blocks} calcule la fermeture des blocs atteignables après un
\code{await} ; la seule transition de phase permise est \code{Sync} →
\code{Deferred}, jamais l'inverse. Décision annexe cruciale : les \emph{IIFE}
(fonctions immédiatement invoquées, \code{(async () => \{ \ldots \})()})
s'exécutent \textbf{à leur site d'appel} --- sans cette règle, aucune ligne de
relation n'était jamais produite pour ce motif très courant en React
(\code{useEffect} avec une IIFE asynchrone interne).

\textbf{Refusé.} Un champ de bloc pour représenter l'\code{await}, plutôt
qu'une variante d'arête : \code{BasicBlock} et \code{CFG} sont construits en
plus de deux cents endroits du code ; une variante d'arête est
\emph{additive} (tout \code{match} existant sur \code{EdgeKind} a déjà un bras
\code{_}), un champ de bloc aurait exigé de revoir chacun des deux cents
sites.

\textbf{Conséquences.} Aucune régression mesurée ; le fait \og la suite d'un
\code{await} est différée\fg{} devient disponible aux relations en aval
(\code{writers}, \code{registrations}).

\textbf{ADR liés.} §1 (l'endroit exact où la scission a lieu) amendé par
\adr{36} §6 : l'expression attendue s'évalue \emph{avant} la suspension, pas
après.

\textbf{Fichiers.} \file{src/ir/cfg.rs} (\code{EdgeKind::Await}).

\textbf{Chapitre.} \cref{chap:ir}.

% ===========================================================================
\section{ADR-036 --- Relation d'appels}
\label{adr:36}
\noindent\textit{2 septembre 2026. Statut : Accepted.}

\textbf{Contexte.} Le corps d'un effet n'était lisible par les règles qu'à
travers trois relations existantes (écriture, lecture, enregistrement) :
aucune ne disait simplement \og cette fonction a été appelée ici\fg{}.

\textbf{Décision.} Un second \textbf{canal} de la \emph{même} marche
(\og the setter walk's second output, not a second walk\fg{}) : une ligne
\code{(name, receiver, phase)} par appel, \emph{sans} les valeurs des
arguments (question distincte, laissée à l'issue \issue{67}). Le nom de
l'appelé est une garde \textbf{obligatoire} : c'est la première relation du
catalogue non bornée en cardinalité, et un appel sans nom résolu est rejeté
même s'il est caché dans un combinateur \code{any_of}. Plafond \Warning{}.
Ancre \code{render_calls} ; quantificateur \code{none} (une ligne
\emph{manquante} est la direction non sound, donc la direction sûre à
choisir : mieux vaut une règle qui tire à tort qu'une qui se tait à tort).
Option \code{elements: component | host | any} ; ancre \code{elements}.

\textbf{Refusé.} Exposer les valeurs des arguments dans cette même relation :
une question à part (\issue{67}), qui mériterait sa propre relation plutôt
qu'un élargissement de celle-ci.

\textbf{Conséquences.} Les règles de \og calls\fg{} (audit d'API interdite,
politique d'équipe) deviennent exprimables sans marcher le CFG elles-mêmes.

\textbf{ADR liés.} Prépare le terrain pour \adr{35} §1 (la position exacte de
la scission \code{await}, amendée par cet ADR §6).

\textbf{Fichiers.} \file{src/engine/setters.rs} (second canal de la même
marche).

\textbf{Chapitre.} \cref{chap:relations}.

% ===========================================================================
\section{ADR-037 --- Relation de lecture de slot}
\label{adr:37}
\noindent\textit{2 septembre 2026. Statut : Accepted.}

\textbf{Contexte.} Toute la relation d'écriture (\relation{slot_writers})
n'avait pas de miroir côté lecture : aucune règle ne pouvait exprimer
\og ce slot est lu ici\fg{}.

\textbf{Décision.} Une image miroir exacte de \relation{slot_writers} :
mêmes colonnes \code{region} et \code{phase}. Troisième canal de la même
marche unique. \code{reads_in} ne traverse \emph{pas} un \code{FnLit}
imbriqué (sinon $\Top$ serait compté deux fois, une fois côté lecture et une
fois côté appelant). La phase fait partie de l'\emph{identité} de la ligne --
sinon une ligne \code{Cleanup} distincte d'une ligne \code{Effect}
disparaîtrait par fusion accidentelle. La marche devient une traversée
\og réelle\fg{} (et non un simple accumulateur) dès qu'un canal quelconque
est actif.

\textbf{Refusé --- reporté, pas tranché.} La moitié \og la lecture
atteint-elle le rendu ?\fg{} : une question de \emph{taint} (propagation de
provenance) distincte, laissée à \adr{41} (dépendance de rendu).

\begin{piege}
La même cécité que celle du côté écriture existe côté lecture : \og the same
blindness on the setter side is an FN\fg{}, déposée comme issue \issue{130}.
Ne pas supposer que corriger un côté de la marche corrige l'autre.
\end{piege}

\textbf{Conséquences.} La lecture d'un slot devient une relation de première
classe, symétrique à l'écriture.

\textbf{ADR liés.} La moitié \og taint\fg{} refusée ici est reprise par
\adr{41}.

\textbf{Fichiers.} \file{src/engine/setters.rs} (troisième canal).

\textbf{Chapitre.} \cref{chap:relations}.

% ===========================================================================
\section{ADR-038 --- Une écriture est une écriture où qu'elle soit écrite}
\label{adr:38}
\noindent\textit{2 septembre 2026. Statut : Accepted, §5 superseded by \adr{40}.}

\textbf{Contexte.} \code{wrap(setN(1))} passait \emph{silencieusement} à
travers la marche d'écriture, alors qu'un \code{setN(1)} nu au même endroit
déclenchait \Error{} \regle{setter-in-render} : un faux négatif de pure
position syntaxique --- le setter était bien appelé, seulement enveloppé dans
un appel de fonction.

\textbf{Décision.} Une traversée unique, \emph{sans porte} conditionnelle sur
la forme syntaxique de l'appel. \regle{setter-in-render} lit enfin la phase
\emph{calculée} par le moteur (\code{Handler}/\code{Deferred} valent preuve de
non-appartenance au rendu) plutôt que la position lexicale. \code{Deferred}
est séparé de $\Top$ dans \code{SetterCallPhase}. Deux conditions must
avant de certifier une écriture en rendu : la phase est \code{Sync}, \emph{et}
la variable appelée \textbf{est} le setter ou une fermeture qui l'appelle
(\code{must_write}). Une seule orthographe retenue pour l'identité d'un
composant (le \emph{display name}) --- ce choix particulier sera remplacé par
\adr{40}.

\begin{exemple}{Le FN de position corrigé}{ex-wrap-setn-adr038}
\begin{lstlisting}[language=TypeScript,caption={Setter enveloppé, appelé en rendu}]
function wrap(f: () => void) {
  f();
}

function Counter() {
  const [n, setN] = useState(0);
  wrap(() => setN(n + 1));
  return <span>{n}</span>;
}
\end{lstlisting}
Sortie observée (\code{\snapshotCommit}) :
\begin{sortie}
  Counter  (1 hooks)  wrapped-render-setter.tsx
    error  setter-in-render  [hook:0]  (line 4:2)  setter `setN` called directly in the render body, move this call into a useEffect or an event handler
       (1 trace step(s), rerun with --trace)

⚠  1 error(s) across 1 file(s).
\end{sortie}
Avant \adr{038}, cette forme --- setter passé en fermeture à une fonction
d'ordre supérieur --- n'avait \og no row in the writer relation at all\fg{}
(citation de l'ADR).
\end{exemple}

\textbf{Refusé.} Rien d'alternatif à la suppression de la porte : la porte
elle-même était le défaut.

\textbf{Conséquences.} Mesure sur corpus : 6\,322 $\to$ 6\,340 findings sur
34\,730 fichiers (27 nouveaux \Warning{}, 9 retraits qui n'étaient pas de
vrais findings mais des doublons de comptage). Sans précalcul de la marche, la
règle aurait tourné 43 fois plus lentement sur le corpus \code{ai-chatbot} ---
justifiant que la relation soit calculée une fois à convergence plutôt qu'à
la demande de chaque règle.

\textbf{ADR liés.} §5 (l'identité de composant par display name) superseded
by \adr{40}.

\textbf{Fichiers.} \file{src/engine/setters.rs} (\code{must_write},
\code{SetterCallPhase}).

\textbf{Chapitre.} \cref{chap:relations}.

% ===========================================================================
\section{ADR-039 --- Une liaison synthétique a une position réelle}
\label{adr:39}
\noindent\textit{2 septembre 2026. Statut : Accepted.}

\textbf{Contexte.} Mesuré : 82 findings sur 7\,146 (1{,}1\,\%) n'avaient
\emph{aucune} position affichable --- le lowering crée des liaisons
synthétiques (des variables temporaires introduites par le désucrage) qui
n'existent pas dans le texte source.

\textbf{Décision.} Chaque instruction synthétique prend la position de
l'expression qu'elle lie. Ce que le code source ne nomme nulle part (par
exemple, un paramètre introduit par un splice) prend le \textbf{site d'appel}
du splice qui l'a introduit. Une marche de relation porte un \code{witness}
\emph{hérité} par défaut, jamais en surcharge d'un témoin déjà présent. Un
finding sans \emph{range} du tout prend la première position de sa chaîne de
témoins --- calculé une seule fois, dans le registre partagé, pour toutes les
règles à la fois.

\textbf{Refusé.} Ajouter un \emph{span} à \code{Terminator::Return}
(reprenant la question laissée ouverte par \adr{25}) : \og not worth it for
what it buys\fg{} au regard du nombre de sites de construction à mettre à
jour.

\textbf{Conséquences.} 82 $\to$ 0 finding sans position ; l'ensemble des
findings lui-même reste inchangé (seule l'attribution de position change).

\textbf{ADR liés.} Referme, sans le résoudre entièrement (\issue{140} reste
ouverte), le même point qu'\adr{25} avait laissé de côté.

\textbf{Fichiers.} Le point de splice (\cref{chap:splice-inlining}) et
\file{src/rules/api/witness.rs} (le registre de positions par défaut).

\textbf{Chapitre.} \cref{chap:splice-inlining}.

% ===========================================================================
\section{ADR-040 --- Identité de composant = \code{ComponentId}}
\label{adr:40}
\noindent\textit{5 septembre 2026. Statut : Accepted ; supersedes \adr{38} §5.}

\textbf{Contexte.} Trois représentations coexistaient pour \og le même\fg{}
composant : une clé \code{(PathBuf, Symbol)}, un \emph{display name} qui
dépendait du \emph{contenu} du fichier (deux fichiers définissant chacun un
\code{Widget} produisaient des collisions), et un \code{Symbol} nu venant de
\code{CompApp}. Deux défauts se terminaient tous deux par un rassurant
\og ✓ … no issues found\fg{} \emph{à tort} : un premier appariement trié par
hasard, et des chemins de fichiers orthographiés différemment selon l'OS.

\textbf{Décision.} Une identité \textbf{interne}, indépendante du contenu :

\rustsnippet{src/ir/component_id.rs}{23}{29}{\code{ComponentId} --- interné, comparable en une instruction}

Le \emph{display name} n'est plus qu'un \textbf{rendu}, calculé seulement à
l'affichage, jamais utilisé comme clé. \code{CompApp.origin} porte le fichier
et le nom exporté. \code{resolve_child} répond l'une de trois réponses, dont
\code{Ambiguous} explicite (plutôt qu'un choix arbitraire silencieux). Tous
les chemins de fichiers du dépôt sont normalisés à la même forme.

\begin{piege}
Le même motif structural revient deux fois dans l'histoire du projet : une
identité \emph{interne}, comparable et indépendante du contenu, à côté d'un
\emph{rendu} calculé seulement à l'affichage. \code{FileId}/\adr{19} l'avait
déjà fait pour les fichiers ; \code{ComponentId}/\adr{40} le refait pour les
composants. Chaque fois qu'une chaîne de caractères sert de \emph{clé} plutôt
que d'étiquette d'affichage, un bug d'identité guette (fil rouge du
\cref{chap:histoire-doctrine}, \og identité $\ne$ rendu\fg{}).
\end{piege}

\textbf{Refusé --- mesuré.} Analyser chaque candidat d'un nom de composant
ambigu, plutôt que de répondre \code{Ambiguous} : 1\,347 références
concernées sur le corpus, une forme quadratique en nombre de composants
(rappel de la contrainte de performance déjà posée par \issue{86} à
\adr{18}).

\textbf{Conséquences.} Le \emph{digest} du corpus (l'empreinte agrégée des
findings) reste \textbf{identique} sur 35\,541 fichiers pour le seul
changement d'identité --- preuve que le renommage n'a, par lui-même, rien
changé au comportement observable. La résolution de composants ambigus passe
de 1\,317 à 1\,348 lignes résolues (29 retraits, 60 ajouts). Commit
\code{806d114} : 107 fichiers, +2\,959/$-$1\,454 lignes ; temps de corpus
873\,s contre 858\,s avant.

\textbf{ADR liés.} Supersedes \adr{38} §5. L'ADR se déclare lui-même
\og Implements: \#7\fg{} et son texte reprend le triple défaut d'identité
que l'issue décrivait ; le tracker, à \code{\snapshotCommit}, marque pourtant
toujours \issue{7} ouverte --- un oubli de fermeture, pas un doute sur la
correction : l'ADR corrige bien ce que l'issue décrivait.

\textbf{Fichiers.} \file{src/ir/component_id.rs}.

\textbf{Chapitre.} \cref{chap:inter-composants}.

% ===========================================================================
\section{ADR-041 --- Dépendance de rendu}
\label{adr:41}
\noindent\textit{24 septembre 2026. Statut : Accepted ; amende \adr{22} §4.}

\textbf{Contexte.} Quelles parties de la sortie JSX d'un composant dépendent
réellement d'un slot d'état, et lesquelles ne font que le \emph{transmettre}
sans jamais l'utiliser pour décider quoi rendre ? \code{Stability}
(\cref{chap:statevalue-stabilite}) ne répond pas à cette question : elle est
booléenne (stable/instable), alors qu'un \code{text.length} ou une allocation
imbriquée demandent une notion plus fine.

\textbf{Décision.} Une analyse \textbf{avant} \emph{séparée}
(\file{engine/render_deps.rs}), sur un composant déjà convergé, calculant des
ensembles \emph{may} de \code{Source} (l'origine d'une valeur) ; distinction
\code{genuine} (usage réel dans une décision de rendu) contre \code{sites}
(propagation pure). Composition le long de l'arbre d'éléments \emph{d'origine
prouvée} : \og absence of a use is a proof, so every unknown is a use\fg{}
--- tout ce que l'analyse ne peut pas exclure compte comme un usage, jamais
l'inverse. Deux règles \Warning{} : \regle{state-lifted-too-high} et
\regle{wasted-subtree-render}. Des options \emph{typées} sur les règles
natives (\code{Rule::options()}, \code{--rule-option}).

\begin{react}[fréquence de déclenchement : continue ou discrète]
Un événement \code{onMouseMove} déclenche à haute fréquence (continu) ; un
\code{onClick} déclenche rarement (discret). \code{Deps::gated} enregistre,
pour une valeur fonction, les captures que son corps n'atteint que sous un
test de l'argument événementiel (issue \issue{148}) --- un événement dont
chaque écriture n'est atteinte qu'à travers un tel test est discret. Un
échec de cette détection classe par défaut en discret : cela ne change jamais
un statut must en may, seulement le classement d'un \Warning{}.
\end{react}

\textbf{Refusé.} Un champ de plus sur le domaine des valeurs
(\code{StateValue}) : cela changerait \emph{toute} comparaison de valeurs,
y compris la clé du cache d'inlining (\code{ComponentCache}) --- un coût de
performance et de complexité disproportionné pour une propriété qui n'est
utile qu'à deux règles.

\textbf{Conséquences.} Un provider de contexte \emph{prouvé} est suivi, pas un
simple usage (issue \issue{145}) : sa valeur atteint, comme
\code{Context(id)}, tout élément imbriqué --- un consommateur peut donc
gaspiller son sous-arbre, et un état peut être \og chez lui\fg{} juste
au-dessus de son provider. Un handler qui écrit un binding de module à côté
d'un état (issue \issue{147}) est suivi jusqu'aux lecteurs de ce binding.
Mesure sur corpus : la baseline passe de 1\,346 à 1\,500 (commit
\code{ba09a62}), puis se stabilise à 1\,493, 1\,494 après corrections de
précision.

\textbf{ADR liés.} Amende \adr{22} §4 (portée des options de règles natives).

\textbf{Fichiers.} \file{src/engine/render_deps.rs}.

\textbf{Chapitre.} \cref{chap:inter-composants} (\cref{def:render-dep}) et
\cref{chap:regles-rendu-rsc} (les deux règles).

% ===========================================================================
\section{ADR-042 --- Les relations sont des produits du moteur}
\label{adr:42}
\noindent\textit{26--27 septembre 2026. Statut : Accepted.}

\textbf{Contexte.} Une revue externe a repéré que le graphe de churn vivait
encore dans \file{rules/helpers/} avec sa \emph{propre} marche de setters,
indépendante de \relation{slot_writers} et sans classification de phase
--- exactement la situation qu'\adr{27} §1 avait commencé à corriger pour les
écritures, mais jamais généralisée. L'issue \issue{26} est le symptôme direct
de cette dérive.

\textbf{Décision.} Une frontière normative, tenue par un \textbf{cliquet}
(un test qui casse la compilation ou la CI si elle est franchie) : \og une
règle lit des lignes et appelle des must-primitives, elle ne marche aucune
syntaxe\fg{}.

\rustsnippet{src/engine/program_relations.rs}{21}{44}{\code{ProgramRelations} --- le graphe de churn construit une fois, en aval du point fixe}

\code{SlotWriter} gagne trois colonnes : \code{owner}, \code{block},
\code{written} (déjà vues aux fiches \adr{30} et \adr{28}, désormais
formalisées ici) :

\rustsnippet{src/engine/setters.rs}{770}{786}{Les colonnes \code{owner}, \code{block}, \code{written} de \code{SlotWriter}}

et \code{Written} lui-même porte la fraîcheur et la valeur écrite :

\rustsnippet{src/engine/written.rs}{37}{58}{\code{Freshness} et \code{Written} --- ce qu'une écriture stocke, pour les preuves de churn}

Une relation \relation{effect_triggers} nouvelle (\textbf{identité}, pas
dépendance de rendu --- délibérément \emph{non} unifiée avec
\relation{render_deps} d'\adr{41}) :

\rustsnippet{src/engine/triggers.rs}{33}{44}{\code{EffectTrigger} --- une ligne par dépendance d'effet qu'un slot fait bouger}

\file{engine/churn.rs} devient un simple \emph{pli} sur deux relations,
combinées par la table de phases (\cref{sec:histoire-doctrine-relations} pour
le détail de cette table) ; les preuves suivent désormais les données :
\og sur tous les chemins\fg{} vit dans \file{engine/dominance.rs}, les gardes
dans \file{engine/guards.rs}. Un catalogue \file{docs/relations.md} documente
toutes les relations publiées.

\textbf{Refusé.} Unifier \relation{effect_triggers} avec \relation{render_deps} :
\og identité $\ne$ dépendance\fg{} --- le \emph{must-rerun} d'un effet a besoin
de savoir \emph{quelle} dépendance a bougé (identité), pas seulement
\emph{si} le rendu en dépend. Dans un amendement du 27 septembre sur les
sites de convergence, un \emph{plus grand} point fixe sur les sites de
revivification est refusé : il \og would read mutual revival as
convergence\fg{} --- deux effets qui se relancent mutuellement à l'infini
semblent \og converger\fg{} vus par le plus grand point fixe, alors qu'ils ne
convergent pas du tout. Le point fixe retenu est le \textbf{plus petit}
(\cref{def:point-fixe}).

\textbf{Conséquences.} \file{rules/helpers/churn.rs} et
\file{rules/helpers/churn_graph.rs} sont supprimés ; \file{engine/churn.rs},
\file{engine/guards.rs}, \file{engine/triggers.rs} et
\file{engine/program_relations.rs} existent désormais ; \file{rules/api/cache.rs}
a disparu (vérifié : absent à \code{\snapshotCommit}). L'issue \issue{26} est
fermée \emph{par construction} (plus de seconde marche à faire dériver), avec
un test de régression dans \file{tests/effect_cycles.rs}. Séquence livrée en
cinq commits séparés, chacun mesuré sur corpus, l'arbre restant vert entre
chaque étape : refactor sans changement de comportement, ajout de
\code{owner}/\code{block}/\code{written}, ajout d'\relation{effect_triggers},
migration de \file{engine/churn.rs} avec re-mesure, puis le test-cliquet et
\file{docs/relations.md}. Faiblesse assumée, héritée d'\adr{29} : les lignes
de cycle n'existent que pour les cycles \emph{vus} par le graphe --- les
arêtes médiées par une prop auraient besoin de faire remonter le slot parent
(\issue{20}, encore ouverte).

\textbf{ADR liés.} Amende \adr{18} (construction et condition d'écrivain
unique), \adr{29} §1 (les lignes viennent désormais du moteur), \adr{21} §4
(\code{ProgramCache} $\to$ \code{ProgramRelations}).

\textbf{Fichiers.} \file{src/engine/program_relations.rs} ;
\file{src/engine/churn.rs} ; \file{src/engine/guards.rs} ;
\file{src/engine/triggers.rs} ; \file{src/engine/setters.rs} ;
\file{src/engine/written.rs}.

\textbf{Chapitre.} \cref{chap:relations} et \cref{chap:churn}.

% ===========================================================================
\begin{resume}
Quarante-deux ADR, quarante-et-une décisions en vigueur et un seul remplacement
complet (\adr{8}, superseded by \adr{15}). Trois formes reviennent tout au
long du catalogue et valent d'être retenues au-delà de la fiche qui les
introduit chacune :

\begin{itemize}
  \item \textbf{Une relation est promue, jamais dupliquée.} À partir
    d'\adr{27}, chaque nouvelle relation (\relation{slot_seeds},
    \relation{context_consumers}, \relation{calls}, \relation{reads},
    \relation{registrations}, le graphe de churn et \relation{effect_triggers})
    est calculée \emph{une fois}, par le moteur, sur la \emph{même} marche
    qu'une relation déjà existante --- jamais par un second parcours du CFG.
    \adr{42} donne la formule qui clôt cette série : \og a fact computed in
    two places is two facts that drift\fg{}.
  \item \textbf{Un non-changement est une décision, pas un oubli.} \adr{20}
    (onze fois) et la table dispersée du \cref{chap:histoire-doctrine}
    (trente-six fois) documentent des simplifications apparentes refusées
    parce qu'elles coûtaient un faux négatif mesurable. Avant de
    \og nettoyer\fg{} une duplication dans ce code, chercher d'abord si un ADR
    l'a déjà justifiée.
  \item \textbf{L'identité n'est jamais un rendu.} \code{FileId} (\adr{19}),
    \code{ComponentId} (\adr{40}), la cellule canonique \code{ContextId}
    (\adr{32}) répètent le même motif : une clé interne, comparable et
    indépendante du contenu, à côté d'une chaîne calculée seulement pour
    l'affichage. Chaque défaut d'identité rencontré dans l'histoire du projet
    venait d'une chaîne de caractères utilisée comme clé.
\end{itemize}

Cette annexe est un catalogue de référence, pas un récit : pour l'histoire
suivie --- pourquoi ces décisions sont venues dans cet ordre, quels fils
rouges les traversent, ce que la mesure sur corpus a changé à la méthode ---
voir \cref{chap:histoire-doctrine}. Pour une décision citée depuis un
chapitre technique, la macro \code{\textbackslash adr\{NN\}} renvoie
directement à la fiche correspondante de cette annexe.
\end{resume}

\section*{Exercices}
\addcontentsline{toc}{section}{Exercices}

\begin{exercice}{Direction d'un refus}{fiches-adr-direction}
Pour cinq ADR de votre choix parmi \adr{9}, \adr{17}, \adr{28}, \adr{34} et
\adr{42}, identifiez dans le champ \og Refusé\fg{} de la fiche correspondante
si l'alternative écartée aurait \emph{ajouté} des findings, en aurait
\emph{retiré}, ou n'aurait rien changé au compte total. Pour chacune, écrivez
la phrase unique qui justifie le refus.
\end{exercice}

\begin{exercice}{Reconstituer une chaîne de remplacement}{fiches-adr-chaine}
En partant d'\adr{8}, suivez la chaîne \og supersedes\fg{}/\og amends\fg{}
jusqu'à la décision en vigueur à \code{\snapshotCommit} pour le domaine des
valeurs d'état. Combien d'ADR faut-il traverser ? Faites de même en partant
d'\adr{11} pour la provenance des diagnostics.
\end{exercice}

\begin{exercice}{Une écriture qui échappe encore}{fiches-adr-echappe}
\adr{38} corrige le cas \code{wrap(setN(1))}. Construisez, sur le modèle de
l'exemple de cette fiche, une variante à deux niveaux d'enveloppe
(\code{wrap(() => wrap2(() => setN(n + 1)))}) et vérifiez avec
\code{cargo run -q -- check} si elle est toujours détectée. Si oui, quelle
propriété de la décision (\og une traversée, sans porte\fg{}) l'explique ?
\end{exercice}

\begin{exercice}{Écrire un ADR}{fiches-adr-ecrire}
Au format des fiches de cette annexe (Contexte, Décision, Refusé,
Conséquences), rédigez un ADR-043 fictif qui autoriserait le point fixe à
descendre dans les appelés \code{Unknown} d'\adr{9} --- puis un second ADR qui
refuse ce changement. Lequel des deux respecte l'invariant \og faux négatifs
interdits\fg{} de \file{CLAUDE.md} ? Lequel respecte la politique historique
\og FP-averse\fg{} d'\adr{9} ? Les deux réponses sont-elles la même décision ?
\end{exercice}
